\documentclass[11pt]{article}

\usepackage[margin=1in]{geometry}
\usepackage{amsmath}
\usepackage{amssymb}
\usepackage{booktabs}
\usepackage[hidelinks]{hyperref}
\usepackage{url}

\title{vkGSplat: Is 3D Gaussian Splatting a Good Integration Layer for DLSS-Class Reconstruction?}
\author{vkGSplat research note}
\date{Started May 5, 2026}

\begin{document}
\maketitle

\section{Executive Thesis}

The answer is \textbf{qualified yes}. 3D Gaussian Splatting is a promising way
to integrate Monte Carlo denoising, radiance caching, temporal anti-aliasing,
temporal super-resolution, supersampling, frame interpolation, and learned
reconstruction, but only if vkGSplat treats 3DGS as a \emph{shared persistent
reconstruction state}, not as a standalone renderer or a bag of unrelated
tricks.

The central claim to test is:

\begin{quote}
A single explicit 3D/4D Gaussian state can carry geometry, visibility, motion,
radiance, subpixel coverage, variance, and confidence across frames. If that
state can feed multiple reconstruction tasks with less ghosting, less shimmer,
and lower variance than separate image-space systems, then 3DGS is a good
integration layer for vkGSplat.
\end{quote}

The software-stack consequence is important: 3DGS is \emph{not} the whole
renderer, and DLSS-style reconstruction is \emph{not} the acquisition stage.
vkGSplat should keep software renderers as replaceable acquisition producers.
A CPU ray tracer, an Apple Metal backend, a future CUDA software ray tracer, a
\textsc{nvdiffrast}-style CUDA triangle rasterizer, and a GraphDECO/\textsc{gsplat}
Gaussian rasterizer should all be allowed to produce the same seed-buffer
contract. The 3DGS/DLSS-class research layer begins after that boundary: it
consumes color, depth, primitive identity, motion, variance, and auxiliary
signals, then performs history, denoising, super-resolution, frame generation,
and optional neural residual reconstruction.

The core framework therefore does \emph{not} require massive training. The
initial stack is analytic and systems-oriented: software acquisition, seed
buffer extraction, reprojection, depth/identity rejection, SVGF-style
denoising, temporal accumulation, Gaussian state updates, and CUDA tile
scheduling. Training appears only as an optional later residual model. The
research bet is that explicit geometry, motion, variance, identity, and
confidence should reduce what a neural model must learn, not force vkGSplat to
replicate a proprietary DLSS-scale training pipeline.

This does \emph{not} mean vkGSplat can directly clone DLSS 4.5 or the
announced DLSS 5 target. NVIDIA's DLSS 4.5 combines a second-generation
transformer Super Resolution model, Dynamic Multi Frame Generation, 6X Multi
Frame Generation, Reflex, mature engine integration, optical flow, training
data, and heavy production tuning \cite{nvidia2026dlss45,nvidia2026dlssdeveloper}.
NVIDIA has also announced DLSS 5 as a real-time neural rendering system that
uses game color and motion-vector inputs to synthesize photoreal lighting and
material detail while remaining anchored to the game scene \cite{nvidia2026dlss5}.
vkGSplat's stronger research position is different:

\begin{quote}
DLSS reconstructs from images plus engine buffers. vkGSplat can reconstruct from
images plus explicit Gaussian scene, motion, and radiance state.
\end{quote}

\section{Training Requirement}

vkGSplat should be evaluated in three increasingly expensive modes:

\begin{enumerate}
  \item \textbf{No-training core.} The required baseline is deterministic:
        seed acquisition, motion-map derivation, history rejection, temporal
        accumulation, denoising, Gaussian update, and Gaussian projection are
        implemented as explicit CPU/Metal/CUDA kernels. This mode must produce
        useful results before any neural model is introduced.
  \item \textbf{Small calibrated residual.} If the analytic path exposes stable
        Gaussian feature planes, a compact residual network may learn bounded
        corrections for high-frequency shading, specular residuals, or
        difficult disocclusions. This model is optional and should be tested
        against an equal-size image-only model.
  \item \textbf{DLSS-class product model.} A large model trained across many
        games/scenes is out of scope for the core framework. It belongs only to
        a future productization path, not to the proof that the software stack
        and Gaussian state are useful.
\end{enumerate}

This distinction is central. A failure to match DLSS product quality does not
falsify vkGSplat. The falsifiable core question is whether explicit
acquisition-independent Gaussian state improves reconstruction quality,
stability, or cost before large-scale training enters the loop.

\section{What ``Good Integration Layer'' Means}

3DGS is useful here only if it satisfies four criteria:

\begin{enumerate}
  \item \textbf{Shared state:} the same Gaussian state helps more than one
        feature: TAA, super-resolution, denoising, radiance caching, and
        interpolation.
  \item \textbf{Better address space:} history is attached to world, object, or
        path evidence, not only previous-frame pixels.
  \item \textbf{Better uncertainty:} each Gaussian carries confidence, variance,
        visibility age, primitive identity, material identity, and motion
        residuals.
  \item \textbf{Compute-friendly:} update, query, and splat operations map to
        tiled CUDA/Metal-style kernels without turning into a training system.
  \item \textbf{Acquisition-independent:} the layer consumes the same seed
        buffers no matter whether the evidence came from CPU, Metal, CUDA
        software rasterization, CUDA ray traversal, or Gaussian splatting.
\end{enumerate}

The answer becomes \textbf{no} if each feature needs its own incompatible
Gaussian state, if the state ghosts more than image history, or if the cost
beats neither a strong TAAU/SVGF baseline nor a small neural resolver.

The implementation focus is CUDA. CPU code is only a small oracle for
debugging and expected-value checks; the research artifact should answer
whether a shared Gaussian state is a good \emph{GPU} integration layer.

\section{Software-Stack Boundary}

This note is a companion to the broader vkGSplat system paper in
\texttt{papers/vkGSplat.tex}. The system paper argues for a compute-first
Vulkan path for synthetic data generation: capture existing Vulkan/Wicked
workloads, lower acquisition to CUDA-class compute, and preserve a stable
seed-buffer ABI. This note deliberately starts after that boundary. It asks
whether a persistent Gaussian state is the right reconstruction and reuse
layer once seed evidence already exists. Therefore the two papers should
cross-reference each other rather than duplicate material: the system paper
owns the Vulkan/SDG/hardware-allocation argument, while this note owns the
3DGS reconstruction, tensorization, radiance-cache, and frame-generation
algorithms.

The current vkGSplat stack should be read as two independent layers joined by a
stable seed ABI:

\[
\begin{gathered}
\text{Vulkan/Wicked command stream or captured scene state}\\
\downarrow\\
\text{software acquisition renderer}\\
\downarrow\\
\text{seed buffers}\\
\downarrow\\
\text{3DGS/DLSS-class reconstruction layer}
\end{gathered}
\]

The acquisition layer answers the question: ``what did the camera observe this
frame?'' It may use rasterization, ray traversal, shadow/light visibility, or
Gaussian splatting. Its job is to emit structured measurements, not final
beauty frames. A seed frame should contain at least:

\begin{itemize}
  \item linear color or radiance,
  \item depth and NDC depth,
  \item primitive/object/material identity,
  \item camera and previous-camera matrices,
  \item screen-space motion or enough state to derive it,
  \item optional normal, albedo, variance, sample count, and path metadata.
\end{itemize}

The reconstruction layer answers a different question: ``how should these
measurements be reused over space, time, and scale?'' It owns temporal
anti-aliasing, denoising, super-resolution, radiance-cache reuse, frame
interpolation, Gaussian state maintenance, and any learned residual. It should
not branch on whether the seed came from CPU, Metal, or CUDA acquisition.

This separation changes how to use existing CUDA renderers. \textsc{nvdiffrast}
is a source-level reference for a CUDA triangle acquisition backend: triangle
setup, bin rasterization, coarse rasterization, and fine rasterization can
produce depth, primitive ID, barycentrics, and coverage
\cite{laine2020nvdiffrast}. The original GraphDECO CUDA rasterizer released
with 3DGS provides the canonical Gaussian acquisition path
\cite{kerbl2023gs}; \textsc{gsplat} adds a maintained library reference for
packed/unpacked projection, tile intersections, flattened IDs, and
front-to-back compositing \cite{ye2024gsplat}. None of these should leak
private assumptions into the DLSS-class reconstruction layer. They are
producers of the same seed ABI.

The clean test is whether a recorded seed buffer can be replayed through the
reconstruction layer with no acquisition backend present. If it can, the stack
is modular. If reconstruction code requires calling a particular renderer, the
boundary has failed.

\section{Relevant Research}

\subsection{Core 3DGS}

3D Gaussian Splatting introduced explicit anisotropic Gaussian primitives with
fast visibility-aware rendering for real-time radiance-field view synthesis
\cite{kerbl2023gs}. For vkGSplat, the important lesson is not just fast splats.
It is that an explicit primitive can be optimized, projected, sorted, blended,
and inspected.

\subsection{Anti-Aliasing and Scale-Aware Splatting}

Mip-Splatting identifies sampling-rate changes, camera distance, and focal
length as causes of 3DGS aliasing, then combines 3D smoothing with a 2D
Mip/box-filter approximation \cite{yu2024mipsplat}. Multi-Scale 3DGS,
Analytic-Splatting, AAA-Gaussians, LOD-GS, and AA-Splat all strengthen the
case that splats need footprint-aware filtering before they can be a robust
temporal reconstruction layer
\cite{yan2024multiscale,liang2024analytic,steiner2025aaa,yang2025lodgs,suh2026aasplat}.

\subsection{Radiance Caching and Monte Carlo Reconstruction}

GSCache is the closest direct precedent: it uses 3D Gaussian splatting as a
multi-level path-space radiance cache for volume path tracing
\cite{bauer2026gscache}. Neural Radiance Caching and Neural Control Variates
are conceptual ancestors for learned or corrected variance reduction
\cite{muller2021nrc,muller2020ncv}. SVGF is the practical image-space baseline
for one-sample-per-pixel real-time path-traced reconstruction
\cite{schied2017svgf}; KPCN is the classic learned Monte Carlo denoising
baseline using auxiliary buffers \cite{bako2017kpcn}. ReSTIR shows how
spatiotemporal sample reuse and biased/unbiased estimators can coexist in a
real-time renderer \cite{bitterli2020restir}. StochasticSplats and Unified
Gaussian Primitives show that Monte Carlo estimators and Gaussian
representations can be coupled more deeply than normal raster splatting
\cite{kheradmand2025stochasticsplats,zhou2024unified}.

\subsection{Temporal and Motion-Aware Gaussians}

4DGS, Spacetime Gaussian Feature Splatting, 4D-Rotor Gaussian Splatting,
Gaussian-Flow, GaussianFlow, MotionGS, and KF-GS all point toward Gaussians as
motion-carrying state, not only static scene primitives
\cite{wu2024dynamic4dgs,li2023spacetime,duan2024rotor4dgs,lin2023gaussianflow,gao2024gaussianflow,zhu2024motiongs,liu2026kfgs}.
RetimeGS is especially relevant because it targets continuous-time 4DGS
retiming and ghost-free arbitrary-timestamp interpolation
\cite{wang2026retimegs}.

\subsection{DLSS, FSR, and XeSS}

DLSS, FSR, and XeSS set the product bar for temporal upscaling and frame
generation. DLSS 4 uses transformer models for Super Resolution, DLAA, and Ray
Reconstruction; DLSS 4.5 upgrades Super Resolution and adds Dynamic Multi Frame
Generation / 6X MFG on supported hardware
\cite{nvidia2025dlss4,nvidia2025dlss4research,nvidia2026dlss45}. DLSS 5 raises
the comparison bar again: the public target is no longer only upscaling and
frame generation, but a neural renderer that infers lighting/material detail
from game outputs while preserving temporal control \cite{nvidia2026dlss5}.
XeSS is AI-based temporal super sampling using jittered color, motion vectors,
and depth \cite{intel2026xess}. AMD FSR includes temporal upscaling, frame
generation, and newer ray regeneration/radiance-caching branding
\cite{amd2026fsr}.

\section{Unified vkGSplat Architecture}

\subsection{One State, Many Views}

The proposed integration layer is a persistent state:

\[
  G_j =
  (\mu_j, \Sigma_j, R_j, c_j, f_j, L_j, v_j, m_j,
   \tau_j, id_j, \dot{\mu}_j, \nu_j, r_j).
\]

Where:

\begin{itemize}
  \item $\mu_j, \Sigma_j, R_j$ are position, covariance, and orientation,
  \item $c_j$ is reconstructed color or albedo-like appearance,
  \item $f_j$ is a compact learned or hand-authored feature vector,
  \item $L_j$ is radiance/cache content, possibly split by path depth,
  \item $v_j$ is variance,
  \item $m_j$ is sample mass,
  \item $\tau_j$ is age or last-seen time,
  \item $id_j$ stores primitive/material/motion group identity,
  \item $\dot{\mu}_j$ is velocity,
  \item $\nu_j$ is visibility confidence,
  \item $r_j$ is residual error, such as projected flow mismatch.
\end{itemize}

For implementation, split this into a shared core plus task-specific subslots:

\[
  G_j = (C_j, A_j^{taa}, A_j^{sr}, A_j^{rad}, A_j^{mot}, A_j^{nn}).
\]

The shared core is:

\[
  C_j = (\mu_j,\Sigma_j,R_j,\bar{n}_j,id_j,m_j,\tau_j,\nu_j,r_j).
\]

Task subslots are:

\[
  A_j^{taa}=(c_j,v_j),\quad
  A_j^{sr}=(c_j,\Sigma_{screen,j},m_j),\quad
  A_j^{rad}=(L_j^0,\ldots,L_j^K,V_{L,j}^0,\ldots,V_{L,j}^K),
\]

\[
  A_j^{mot}=(\dot{\mu}_j,\mu_j^0,\mu_j^1,R_j^0,R_j^1),
  \quad
  A_j^{nn}=f_j.
\]

This decomposition is the compromise that makes the integration thesis
testable. The address, visibility, covariance, identity, and confidence are
shared. Payloads may differ by feature. If the payloads force incompatible
cores, the unified-state idea fails.

The state exposes multiple projections:

\begin{itemize}
  \item \textbf{TAA projection:} splat stable history into current pixels.
  \item \textbf{Super-resolution projection:} resolve target-resolution pixels
        from accumulated subpixel Gaussian evidence.
  \item \textbf{Radiance-cache query:} query indirect radiance at path hits.
  \item \textbf{Frame-interpolation projection:} evaluate state at $t_\alpha$
        and splat an intermediate frame.
  \item \textbf{Neural feature projection:} render Gaussian confidence,
        variance, radiance, age, and motion residuals as inputs to a small
        learned resolver.
\end{itemize}

\subsection{Pipeline}

Per frame:

\begin{enumerate}
  \item Acquire a low-resolution/low-spp seed frame from any backend that
        satisfies the seed ABI: current CPU fixture, Metal validation path,
        future CUDA ray/raster path, or optional Gaussian raster path.
  \item Extract samples: color, depth, NDC depth, normal, primitive ID,
        material ID, motion vector, path contribution, and optional path depth.
  \item Update the Gaussian state with visibility and confidence gates.
  \item Produce task-specific outputs from the same state:
        TAA history, SR resolve, radiance cache, interpolation state, and
        neural features.
  \item Compare against current-frame samples and references, then decay or
        invalidate stale Gaussians.
\end{enumerate}

\subsection{Why This Could Beat Separate Systems}

Separate systems duplicate partial history:

\begin{itemize}
  \item TAA stores image-space color history.
  \item SVGF stores moments and filtered radiance.
  \item Super-resolution stores subpixel history.
  \item Frame interpolation stores optical-flow history.
  \item Radiance caching stores lighting history.
\end{itemize}

The 3DGS hypothesis is that these are shadows of the same underlying state:
where scene evidence was, how it moved, how visible it is, what radiance it
carried, and how reliable it is. A unified Gaussian state is attractive because
it can make that evidence explicit.

\section{Formal Model and Algorithms}

\subsection{Observations}

At frame $t$, vkGSplat receives a set of observations:

\[
  S_t = \{s_i^t\}_{i=1}^{N_t}.
\]

Each observation can come from a visible pixel, a subpixel jitter sample, or a
path-tracing event:

\[
  s_i^t =
  (x_i, n_i, c_i, L_i^0,\ldots,L_i^K, d_i, z_i, u_i,
   \rho_i, id_i, q_i, \sigma_i^2).
\]

Where $x_i$ is world position, $n_i$ is normal, $c_i$ is color, $L_i^k$ is a
radiance contribution for path-depth level $k$, $d_i$ is depth, $z_i$ is NDC
depth, $u_i$ is screen-space motion, $\rho_i$ is roughness/material metadata,
$id_i$ is primitive/material/motion identity, $q_i$ is sample quality, and
$\sigma_i^2$ is a local variance estimate.

\subsection{Sample Extraction}

For a raster or ray-traced seed pixel $p=(p_x,p_y)$, unproject with the current
inverse view-projection matrix:

\[
  x_i =
  \Pi_t^{-1}(p,z_i)
  =
  \left(VP_t\right)^{-1}
  \begin{bmatrix}
    2(p_x+\frac{1}{2})/W - 1\\
    2(p_y+\frac{1}{2})/H - 1\\
    z_i\\
    1
  \end{bmatrix}.
\]

After homogeneous divide, $x_i$ becomes the world-space address for history.
This matches the existing vkGSplat reprojection contract: depth plus camera
matrices are already enough to derive a current-to-previous motion map. The
Gaussian path simply stores the result in a persistent state.

For temporal super-resolution with scale factor $s\in\{2,4\}$, the low-res
pixel contributes to high-res coordinates:

\[
  P_i = s\,p + s\,\delta_t,
\]

where $\delta_t$ is the subpixel jitter. The sample therefore updates both a
world-space Gaussian and a subpixel coverage estimate.

For path-traced samples, the extraction point is the path event rather than the
pixel center:

\[
  s_i^t =
  (x_{\ell}, n_{\ell}, \omega_o, \omega_i, \beta_{\ell}L_{\ell},
   pdf_{\ell}, \ell, id_{\ell}),
\]

where $\ell$ is path depth and $\beta_{\ell}$ is throughput. The first
radiance-cache prototype should keep only $\ell=1$ rough/diffuse indirect
events so the experiment is hard to fool but still easy to debug.

\subsection{Compatibility Gates}

An observation updates a Gaussian only if it is compatible. Define:

\[
  W_{ij} =
  w_x\,w_n\,w_d\,w_{id}\,w_\rho\,w_v\,w_r\,w_q.
\]

Spatial compatibility:

\[
  w_x =
  \exp\left(-\frac{1}{2}(x_i-\mu_j)^T\Sigma_j^{-1}(x_i-\mu_j)\right).
\]

Normal compatibility:

\[
  w_n =
  \max(0, n_i \cdot \bar{n}_j)^{\gamma_n}.
\]

Depth compatibility after projection:

\[
  w_d =
  \begin{cases}
  1, & |\hat{d}_{ij} - d_i| < \epsilon_d,\\
  0, & \text{otherwise}.
  \end{cases}
\]

Identity compatibility:

\[
  w_{id} =
  \begin{cases}
  1, & id_i = id_j,\\
  \eta_{id}, & id_i \neq id_j \text{ but material-compatible},\\
  0, & \text{otherwise}.
  \end{cases}
\]

Motion compatibility:

\[
  w_v =
  \exp\left(-\frac{\|u_i - \hat{u}_j\|^2}{2\sigma_u^2}\right).
\]

Material/roughness compatibility:

\[
  w_\rho =
  \exp\left(-\frac{|\rho_i-\rho_j|^2}{2\sigma_\rho^2}\right).
\]

Flow residual compatibility:

\[
  w_r =
  \exp\left(-\frac{r_j^2}{2\sigma_r^2}\right).
\]

Sample-quality compatibility:

\[
  w_q = \mathrm{sat}(q_i).
\]

These gates are the anti-ghosting core. A prototype may begin with only
$w_x,w_d,w_{id}$, then add normals, motion, and residuals as the scenes become
harder.

\subsection{Assignment}

For each observation, find the best compatible Gaussian:

\[
  j^* = \arg\max_j W_{ij}.
\]

If $W_{ij^*} < \epsilon_W$, spawn a new Gaussian. Otherwise update $G_{j^*}$.
The CUDA prototype should use a hashed 3D grid plus primitive buckets,
fixed-radius cell neighborhoods, and bounded candidate lists. A scalar CPU path
may exist only as a tiny correctness oracle; it is not the implementation
target.

\subsection{State Update}

Let $\alpha_{ij}$ be the update weight:

\[
  \alpha_{ij} =
  \frac{W_{ij}}{m_j + W_{ij} + \epsilon}.
\]

Mean update:

\[
  \mu_j' = (1-\alpha_{ij})\mu_j + \alpha_{ij}x_i.
\]

Color update:

\[
  c_j' = (1-\alpha_{ij})c_j + \alpha_{ij}c_i.
\]

Radiance level update:

\[
  (L_j^k)' = (1-\alpha_{ij})L_j^k + \alpha_{ij}L_i^k.
\]

Velocity update:

\[
  \dot{\mu}_j' =
  (1-\alpha_{ij})\dot{\mu}_j
  + \alpha_{ij}\frac{x_i-\mu_j^{prev}}{\Delta t}.
\]

Sample mass and age:

\[
  m_j' = \lambda_m m_j + W_{ij},
  \qquad
  \tau_j' = t.
\]

Confidence combines mass, age, variance, and residual consistency:

\[
  \nu_j =
  \mathrm{sat}\left(
    \frac{m_j}{m_j + m_0}
    \exp(-\lambda_\tau(t-\tau_j))
    \frac{1}{1+v_j}
    \exp(-r_j^2 / 2\sigma_r^2)
  \right).
\]

Variance can be updated with an exponential moment:

\[
  v_j' =
  (1-\alpha_{ij})v_j
  + \alpha_{ij}\|c_i-c_j\|^2.
\]

\subsection{Covariance Update}

For a surface-attached Gaussian, use a tangent-frame covariance:

\[
  \Sigma_j =
  R_j
  \begin{bmatrix}
    \sigma_{t1}^2 & 0 & 0\\
    0 & \sigma_{t2}^2 & 0\\
    0 & 0 & \sigma_n^2
  \end{bmatrix}
  R_j^T.
\]

The normal-axis variance should stay small:

\[
  \sigma_n^2 = \min(\sigma_n^2, \epsilon_n^2).
\]

For temporal super-resolution, tangent variances should track projected pixel
footprint and sample density:

\[
  \sigma_t^2 =
  \max(\sigma_{min}^2,
       \beta_f A_{pixel\rightarrow world}
       /(m_j+\epsilon)).
\]

This makes Gaussians shrink as evidence accumulates but prevents unstable
point-like splats.

\subsection{Projection to Pixels}

Project each Gaussian into camera $C_t$:

\[
  p_j = \Pi_t(\mu_j).
\]

Approximate the screen-space covariance with the projection Jacobian $J_j$:

\[
  \Sigma_{screen,j} = J_j\Sigma_jJ_j^T + \Sigma_{pixel}.
\]

Pixel contribution:

\[
  a_{pj} =
  \nu_j
  \exp\left(
    -\frac{1}{2}(p-p_j)^T
    \Sigma_{screen,j}^{-1}
    (p-p_j)
  \right).
\]

Visibility/depth testing gates $a_{pj}$ before blending. The first CUDA
prototype can use per-tile depth gates and bounded front-to-back candidate
lists; later versions can use sorted or weighted visibility like ordinary 3DGS.

\subsection{TAA Resolve}

The Gaussian history color for pixel $p$ is:

\[
  H_t(p) =
  \frac{\sum_j a_{pj}c_j}{\sum_j a_{pj}+\epsilon}.
\]

Blend with the current seed color $C_t(p)$:

\[
  O_t(p) =
  (1-\beta_p)C_t(p) + \beta_p H_t(p).
\]

Where:

\[
  \beta_p =
  \mathrm{sat}\left(
  \beta_0
  \frac{\sum_j a_{pj}}{\sum_j a_{pj}+a_0}
  \frac{1}{1+\bar{v}_p}
  \right).
\]

This is the direct replacement for image-space history. It should fail closed:
if confidence is low, output the current seed.

\subsection{Temporal Super-Resolution Resolve}

For output pixel $P$ at high resolution and scale factor $s$, use the same
projection equation but evaluate at the high-resolution pixel center:

\[
  O^{SR}_t(P) =
  \frac{\sum_j a_{Pj}c_j + \lambda C^{up}_t(P)}
       {\sum_j a_{Pj}+\lambda}.
\]

$C^{up}_t$ is a simple upsampled current frame. $\lambda$ prevents stale
history from dominating disoccluded regions. This is the first test of
integration: the TAA and SR resolves must use the same $G_j$, not two separate
caches.

For 2x and 4x experiments, the only scale-specific term should be the pixel
footprint:

\[
  \Sigma_{pixel}^{SR}(s) =
  \frac{1}{s^2}\Sigma_{pixel}.
\]

If the method needs a different Gaussian state for 2x and 4x, it fails the
integration test. A single state may still expose different projection
footprints.

\subsection{Radiance Cache Query}

At a shading point $x$ with normal $n$, outgoing direction $\omega_o$, roughness
$\rho$, and path-depth level $k$, query:

\[
  \hat{L}^k(x,\omega_o) =
  \frac{\sum_j W_j(x,n,\omega_o,\rho,k)L_j^k}
       {\sum_j W_j(x,n,\omega_o,\rho,k)+\epsilon}.
\]

The query weight reuses the same gates as TAA plus roughness/path-depth
compatibility. The first experiment should query only $k=1$ diffuse or rough
indirect radiance. Direct lighting remains sampled normally.

Radiance cache update uses the same assignment but separate moments per path
level:

\[
  (M_{L,j}^k)' =
  (1-\alpha_{ij})M_{L,j}^k
  + \alpha_{ij}L_i^k,
\]

\[
  (V_{L,j}^k)' =
  (1-\alpha_{ij})V_{L,j}^k
  + \alpha_{ij}\|L_i^k - M_{L,j}^k\|^2.
\]

When used as a biased indirect-light estimate, the cache contributes directly.
When used as a safer denoising feature, it is projected into a feature frame:

\[
  F_L(p,k) =
  \frac{\sum_j a_{pj}L_j^k}{\sum_j a_{pj}+\epsilon}.
\]

The unbiased/control-variate path is a later experiment:

\[
  \hat{I} =
  G_\theta +
  \frac{1}{N}\sum_i \frac{f(x_i)-g_\theta(x_i)}{p(x_i)}.
\]

The first CUDA prototype should not start here. First prove that the biased
cache reduces variance without unacceptable light leaks.

\subsection{Path-Space Gaussian Reconstruction}

A full path-space Gaussian would live in too many dimensions:

\[
  \xi =
  (x_0,\omega_0,x_1,\omega_1,\ldots,x_\ell,\omega_\ell).
\]

That is not the right first representation for vkGSplat. Instead, use
surface-attached Gaussians as addresses and attach low-dimensional path
subslots:

\[
  A_j^{path}(k,b) =
  (M_{jkb},Q_{jkb},N_{jkb},R_{jkb}),
\]

where $k$ is path depth, $b$ is a coarse directional/material lobe, $M$ is the
mean contribution, $Q$ is the second moment, $N$ is sample count, and $R$ is a
robust residual. A path event chooses a lobe:

\[
  b_i =
  B(\omega_i,\omega_o,\rho_i,material_i).
\]

The updated contribution is:

\[
  Y_i =
  \frac{\beta_i f_i L_i}{p_i},
\]

with throughput $\beta_i$, scattering term $f_i$, incident radiance $L_i$, and
sampling density $p_i$. Use robust clamping only for the biased reconstruction
mode:

\[
  \bar{Y}_i =
  \min(\|Y_i\|, c_\sigma\sqrt{Q_{jkb}})\frac{Y_i}{\|Y_i\|+\epsilon}.
\]

Then update the lobe moments:

\[
  N'_{jkb}=N_{jkb}+w_i,
  \qquad
  M'_{jkb}=M_{jkb}+\frac{w_i}{N'_{jkb}}(\bar{Y}_i-M_{jkb}),
\]

\[
  Q'_{jkb}=Q_{jkb}+\frac{w_i}{N'_{jkb}}(\|\bar{Y}_i\|^2-Q_{jkb}).
\]

This is implementable as a CUDA \texttt{UpdatePathSubslots} kernel.
The first test should use a single diffuse lobe so $b=0$; if that cannot reduce
variance without leaking across primitive boundaries, higher-dimensional
path-space reconstruction is unlikely to rescue the idea.

\subsection{Gaussian Control Variates}

There are two estimator modes, and they should be kept separate in code.

\textbf{Biased reconstruction mode} uses the Gaussian cache as a denoised
estimate:

\[
  O_t(p) =
  C^{direct}_t(p) + \hat{L}^{indirect}_G(p).
\]

This is the fastest useful prototype and should be judged by MSE, temporal
stability, and light-leak masks.

\textbf{Control-variate mode} treats the Gaussian cache as a predictor:

\[
  \hat{I}_{CV} =
  \frac{1}{N}\sum_{i=1}^N
  \left(Y_i - g_G(\xi_i)\right)
  + \hat{G}.
\]

Here $g_G(\xi_i)$ is the Gaussian prediction evaluated at the same path sample
and $\hat{G}$ is the cache's estimate of the corresponding integral. This is
only unbiased if $\hat{G}$ matches the same measure; otherwise it is a biased
residual estimator. The CUDA-facing API should therefore expose:

\begin{verbatim}
RadianceEstimate query_biased_radiance(const ShadingQuery& q) const;
RadianceEstimate query_control_variate(const PathEvent& e) const;
\end{verbatim}

The falsifiable experiment is straightforward: compare equal-spp variance and
bias against a high-spp reference for three modes: no cache, biased Gaussian
cache, and Gaussian residual/control-variate mode. The control-variate path is
worth pursuing only if variance drops faster than bias rises.

\subsection{Gaussian-Guided Sampling}

Once query quality is established, the cache can guide future samples. Use a
mixture proposal:

\[
  p'(y|x) =
  (1-\eta)p_{base}(y|x) + \eta p_G(y|x).
\]

The contribution remains MIS-weighted:

\[
  Y_i' =
  \frac{f(y_i)}{p'(y_i|x_i)}.
\]

For the first vkGSplat prototype, $p_G$ should be very simple: choose one of the
top-$K$ nearby radiance-bearing Gaussians and sample a cosine lobe around its
mean incident direction. The experiment is a kill test. If Gaussian-guided
sampling helps only when the scene is already easy, the cache is useful as a
denoising feature but not as a path-space variance-reduction engine.

\subsection{Motion-Accurate Interpolation}

For midpoint synthesis at $\alpha\in(0,1)$:

\[
  \mu_j^\alpha =
  (1-\alpha)\mu_j^0 + \alpha\mu_j^1.
\]

Constant-velocity fallback:

\[
  \mu_j^\alpha = \mu_j^0 + \alpha\Delta t\,\dot{\mu}_j^0.
\]

Confidence is reduced by forward/backward disagreement:

\[
  \nu_j^\alpha =
  \nu_j^0\nu_j^1
  \exp\left(-\frac{\| \Pi_1(\mu_j^1) -
  (\Pi_0(\mu_j^0)+u_{0\rightarrow1})\|^2}{2\sigma_u^2}\right).
\]

Splat $G_j^\alpha$ into the midpoint camera and compare against a true rendered
middle frame. This is motion-accurate only if silhouettes and depth edges land
correctly, not merely if the image looks plausible.

For rigid objects, the preferred update is not learned velocity but engine
transform:

\[
  \mu_j^\alpha =
  T_{obj}(\alpha)T_{obj}(0)^{-1}\mu_j^0.
\]

For non-linear motion, use a cubic trajectory once four frames are available:

\[
  \mu_j(t) = a_jt^3 + b_jt^2 + c_jt + d_j.
\]

Motion confidence is the minimum of visibility confidence, transform
confidence, and flow-residual confidence. This ensures frame interpolation
fails closed in disoccluded or mismatched regions.

\subsection{Frame Generation and Multi-Frame Generation}

Motion-accurate interpolation generates one or more intermediate frames between
two real frames:

\[
  \alpha_k = \frac{k}{K+1},
  \qquad
  k=1,\ldots,K.
\]

For each $\alpha_k$, evaluate $G_j^{\alpha_k}$ and splat with the same
visibility/confidence gates used by midpoint interpolation. The frame is not
accepted as a confident generated frame unless the per-pixel support is strong:

\[
  S_{\alpha}(p)=\sum_{j\rightarrow p}a_{pj}^{\alpha}\nu_j^{\alpha},
  \qquad
  S_{\alpha}(p) > S_{min}.
\]

Pixels with weak support become holes:

\[
  H_{\alpha}(p)=
  \begin{cases}
  1, & S_{\alpha}(p)\leq S_{min}\ \text{or}\ D_{\alpha}(p)=1,\\
  0, & \text{otherwise}.
  \end{cases}
\]

The generated output is therefore explicitly decomposed:

\[
  O_{\alpha}(p) =
  (1-H_{\alpha}(p))O_{\alpha}^{gauss}(p)
  + H_{\alpha}(p)O_{\alpha}^{fallback}(p).
\]

For a non-neural prototype, $O_{\alpha}^{fallback}$ is nearest real-frame hold
or image-space reprojection. For a neural prototype, it is a bounded residual
resolver conditioned on the hole mask. This is important: vkGSplat should not
pretend 3DGS solved disoccluded background generation. It should make unknown
regions first-class.

For extrapolative frame generation, predict from one real frame:

\[
  \hat{\mu}_j(t+\Delta) =
  \mu_j(t) + \Delta\dot{\mu}_j(t) +
  \frac{1}{2}\Delta^2\ddot{\mu}_j(t).
\]

Uncertainty must grow with prediction horizon:

\[
  \hat{\Sigma}_j(t+\Delta) =
  \Sigma_j(t) + \Delta^2\Sigma_{\dot{\mu},j}
  + \Delta^4\Sigma_{\ddot{\mu},j} + Q_j\Delta,
\]

\[
  \hat{\nu}_j(t+\Delta) =
  \nu_j(t)
  \exp(-\Delta/\tau_{motion})
  \exp(-r_j^2/2\sigma_r^2).
\]

This turns long-horizon frame generation into a confidence problem rather than
a hallucination problem. The first experiment should generate only one
midpoint. The second can generate $K=2$ intermediate frames. Extrapolation
should wait until interpolation passes the silhouette and disocclusion tests.

\subsection{Frame-Generation Experiments}

Use three clips:

\begin{enumerate}
  \item \textbf{Rigid transform clip:} one object with known object transform.
        Expected pass: generated silhouettes match the held-out real frames.
  \item \textbf{Camera strafe clip:} static scene with disocclusion. Expected
        pass: disoccluded holes are masked and do not receive high confidence.
  \item \textbf{Specular/deforming clip:} intentionally hard case. Expected
        pass: residual and confidence metrics predict failure rather than
        hiding it.
\end{enumerate}

Report generated-frame metrics separately from TAA/SR metrics:

\[
  E_{FG} =
  E_{mid} + \lambda_sE_{sil}+\lambda_hE_{hole}
  + \lambda_fE_{flow}.
\]

The shared-state thesis is supported only if frame generation reuses the same
Gaussian core and if its failure masks are useful to the TAA/SR/neural paths.

\subsection{Neural Feature Projection}

The optional neural resolver should not receive an opaque ``Gaussian image.''
It should receive explicit feature planes:

\[
  \Phi_t(p) =
  [H_t(p), F_L(p,0),\ldots,F_L(p,K), \bar{v}_p,
   \bar{\nu}_p, \bar{\tau}_p, \bar{r}_p, D_p, M_p].
\]

Where $D_p$ is a disocclusion likelihood and $M_p$ is a motion-consistency
feature. A compact model can then learn the residual:

\[
  O_t(p) =
  O^{analytic}_t(p) +
  N_\theta(C_t, depth_t, velocity_t, \Phi_t)(p).
\]

The key ablation is equal model size with and without $\Phi_t$. If Gaussian
features do not improve quality or reduce model size, the integration layer is
less useful than hoped.

\subsection{Split, Merge, and Prune}

Split a Gaussian when it is high-confidence but high-variance:

\[
  m_j > m_{split}
  \quad\text{and}\quad
  v_j > v_{split}.
\]

Merge two Gaussians when they are compatible and neither split nor depth gates
object:

\[
  D_{KL}(\mathcal{N}_a,\mathcal{N}_b) < \epsilon_{merge},
  \quad id_a = id_b.
\]

Prune stale Gaussians:

\[
  \nu_j < \epsilon_\nu
  \quad\text{or}\quad
  t-\tau_j > T_{max}.
\]

The first CUDA implementation can skip split/merge and only spawn/prune. Split
and merge become necessary if the shared state blurs thin geometry or grows
without bound.

\subsection{Visibility and Lifetime Update}

Visibility is the part most likely to decide whether the unified state works.
Every frame, a Gaussian can be observed, occluded, disoccluded, or simply
outside the frustum. Treat these as different lifetime events:

\[
  \ell_j^t \in
  \{\mathrm{visible},\mathrm{occluded},\mathrm{disoccluded},
    \mathrm{out\_of\_frustum}\}.
\]

If a projected Gaussian lands behind the current depth buffer by more than a
threshold, it is occluded, not stale:

\[
  \ell_j^t=\mathrm{occluded}
  \quad\text{if}\quad
  d_j^t - D_t(p_j) > \epsilon_{occ}.
\]

If it lands in front of current depth or disagrees with primitive ID, it is
potentially stale:

\[
  \ell_j^t=\mathrm{disoccluded}
  \quad\text{if}\quad
  D_t(p_j)-d_j^t > \epsilon_{front}
  \quad\text{or}\quad
  id_j \neq ID_t(p_j).
\]

Use different confidence decay rates:

\[
  \nu_j^{t+1} =
  \begin{cases}
    \min(1,\nu_j^t+\Delta_\nu), & \ell_j^t=\mathrm{visible},\\
    \lambda_{occ}\nu_j^t, & \ell_j^t=\mathrm{occluded},\\
    \lambda_{dis}\nu_j^t, & \ell_j^t=\mathrm{disoccluded},\\
    \lambda_{frustum}\nu_j^t, & \ell_j^t=\mathrm{out\_of\_frustum}.
  \end{cases}
\]

Require $\lambda_{dis}<\lambda_{occ}$. Occluded objects should be allowed to
survive; disoccluded stale history should disappear quickly. This distinction
is essential for frame interpolation and TAA: otherwise the cache either ghosts
or forgets too aggressively.


\subsection{Tensorized Kernel ABI and M--N--K Estimates}

The intended vkGSplat kernel ABI is tensorized, even when the implementation is
custom CUDA rather than generic PyTorch. The important boundary is:

\begin{quote}
All persistent state, all seed inputs, and all reconstruction outputs should be
represented as contiguous tensor-shaped structure-of-arrays buffers. Irregular
operations such as compaction, grid build, assignment, split/merge, and splat
resolve remain custom CUDA tensor operators.
\end{quote}

This gives vkGSplat two front doors over the same kernels. The research front
door accepts \texttt{torch.Tensor}/\textsc{nvdiffrast} buffers; the runtime
front door accepts Vulkan-exported CUDA memory with the same shapes and
strides. PyTorch should be a caller and allocator for experiments, not the
owner of the renderer architecture.

For the rest of this note, use the following M--N--K convention:

\begin{itemize}
  \item $M$ is the number of observations entering a kernel: raster pixels,
        depth-peel pixels, stochastic samples, or compacted seed hits.
  \item $N$ is the number of active persistent Gaussian records available to
        query, update, project, or resolve.
  \item $K$ is the feature width carried by each observation, Gaussian, or
        projected feature plane. Local nearest-neighbor candidate count is
        written as $C$ so it is not confused with feature width.
\end{itemize}

The acquisition tensor contract should start as:

\[
\begin{aligned}
  R      &\in \mathbb{R}^{B\times H\times W\times 4},
          &&R=(u,v,z/w,id_{tri}+1),\\
  C_{img}&\in \mathbb{R}^{B\times H\times W\times 4},
          &&\text{radiance or color plus confidence/alpha},\\
  X      &\in \mathbb{R}^{B\times H\times W\times 3},
          &&\text{world position from interpolation or ray hit},\\
  Nrm    &\in \mathbb{R}^{B\times H\times W\times 3},
          &&\text{world normal},\\
  V_{px} &\in \mathbb{R}^{B\times H\times W\times 2},
          &&\text{previous-pixel minus current-pixel motion}.
\end{aligned}
\]

A tensorized extraction kernel compacts valid hits into structure-of-arrays:

\[
\begin{aligned}
  X_s       &\in \mathbb{R}^{M_s\times 3}, &
  N_s       &\in \mathbb{R}^{M_s\times 3}, &
  L_s       &\in \mathbb{R}^{M_s\times 4},\\
  p_s       &\in \mathbb{R}^{M_s\times 2}, &
  v_s       &\in \mathbb{R}^{M_s\times 2}, &
  b_s       &\in \mathbb{R}^{M_s\times 2},\\
  q_s       &\in \mathbb{R}^{M_s\times 2}, &
  I_s       &\in \mathbb{Z}^{M_s\times 4}. &
\end{aligned}
\]

Here $q_s=(z/w,\nu)$ and $I_s=(id_{tri},id_{px},id_{batch},flags)$. The
current minimal sample width is therefore

\[
  K_s = 18\text{ fp32 lanes} + 4\text{ int32 lanes}
      \approx 88\text{ bytes per compact sample}.
\]

Material IDs, roughness/metalness, per-sample variance, depth derivatives, and
path-space throughput should be appended as extra tensor columns. The first
CUDA implementation should keep the minimal $K_s$ above, then measure which
extra channels actually improve denoising, temporal stability, or frame
interpolation.

The persistent Gaussian state should also be SoA tensors:

\[
\begin{aligned}
  \mu       &\in \mathbb{R}^{N\times 3}, &
  \Sigma   &\in \mathbb{R}^{N\times 6}, &
  \bar L   &\in \mathbb{R}^{N\times 4},\\
  \sigma_L^2 &\in \mathbb{R}^{N\times 4}, &
  \bar n   &\in \mathbb{R}^{N\times 3}, &
  \bar v   &\in \mathbb{R}^{N\times 3},\\
  h        &\in \mathbb{R}^{N\times K_h}, &
  meta     &\in \mathbb{Z}^{N\times K_i}. &
\end{aligned}
\]

A no-neural analytic state needs roughly $K_g=32$--$48$ fp32 lanes plus
$K_i=4$--$8$ integer lanes per Gaussian: mean, covariance, radiance moments,
normal, velocity, mass/confidence, timestamps, primitive/material identity,
cell key, and flags. At fp32 this is about $160$--$240$ bytes per Gaussian. A
small learned residual can add $K_h=16$--$64$ feature lanes, increasing the
state by $64$--$256$ bytes per Gaussian. This is why the default experiment
should prove value with $K_h=0$ before adding neural features.

Let $P=H W$, $S$ be the number of acquisition samples per pixel, $L$ the number
of depth layers, and $\rho$ the valid-surface hit ratio. Then

\[
  M_{raw}=B P S L,
  \qquad
  M_s\approx \rho M_{raw}.
\]

For DLSS-style super-resolution the output pixel count grows with the linear
scale $r$:

\[
  P_{out}=r^2P,
\]

but $M_s$ does not need to grow with $r$ unless vkGSplat deliberately acquires
extra subpixel samples. This distinction is important: the Gaussian state is
supposed to amortize history and geometry so a $2\times$ or $4\times$ resolve
can increase projected output work without requiring $4\times$ or $16\times$
new shading work every frame.

The first-order $M$ and $N$ budgets are:

\begin{center}
\begin{tabular}{lrrrrr}
\toprule
Mode & $P$ & $M_s$ single layer & $M_s$ 4 jitter/depth & $N$ retained & State memory \\
\midrule
1080p & $2.07$M & $1.2$--$2.0$M & $5$--$8$M & $1.5$--$4$M & $0.25$--$1.0$GB \\
1440p & $3.69$M & $2.2$--$3.5$M & $9$--$14$M & $3$--$7$M & $0.5$--$1.7$GB \\
4K    & $8.29$M & $5$--$8$M   & $20$--$32$M & $6$--$16$M & $1.0$--$4.0$GB \\
\bottomrule
\end{tabular}
\end{center}

The $N$ range assumes that visible surfaces need roughly $0.5$--$1.5$
Gaussians per currently valid pixel after merging, plus a history retention
factor of $1.5$--$3.0$ for temporarily occluded or subpixel evidence. Thin
geometry, particles, foliage, transparencies, and depth peeling push $N$ upward;
flat static walls should merge down aggressively. If $N$ must exceed about
$2P$ for ordinary opaque scenes, the covariance/update rules are probably too
fragmented. If $N<0.25P$ cannot preserve edges and subpixel motion, the state
is probably too coarse for DLSS-class reconstruction.

The key kernel shapes are:

\begin{center}
\begin{tabular}{llll}
\toprule
Kernel & Shape & Main cost & Health target \\
\midrule
Extract hits & $M_{raw}\rightarrow M_s$ & $O(M_{raw})$ & bandwidth-bound, one pass \\
Build grid & $N\rightarrow$ cells & $O(N)$ + scan & $B_{p99}$ bounded \\
Assign samples & $M_s\times C$ & local search & $C\le 32$ typical, $C\le 64$ p99 \\
Reduce updates & $M_s\rightarrow N$ & sort or segmented reduce & fewer atomics than direct update \\
Project splats & $N\rightarrow$ tiles & tile-list emission & reused by TAA/SR/framegen \\
Resolve & $P_{out}\times K_p$ & tile traversal & no duplicated feature passes \\
\bottomrule
\end{tabular}
\end{center}

Projected resolver feature width should begin small:

\[
  K_p \approx 16\text{--}32
  \quad\text{for}\quad
  (L,\alpha,z,n,v,\nu,\sigma^2,id,coverage,age,flags).
\]

where integer IDs can live in parallel integer tensors. A neural resolver may
consume $K_p=32$--$96$ channels, but the pass/fail experiment should compare
against the analytic $K_p\le 32$ case first.

These estimates imply three CUDA milestones:

\begin{enumerate}
  \item \textbf{M0 tensor extraction:} $M_s\le 2$M, $N=0$, $K_s=18+4$.
        Demonstrate \textsc{nvdiffrast}/Vulkan seed tensors can become compact
        sample tensors in less than one millisecond at 1080p-class sizes on a
        desktop CUDA GPU.
  \item \textbf{M1 persistent update:} $M_s=1$--$4$M, $N=1$--$3$M,
        $K_g\le 48$. Demonstrate bounded grid occupancy, stable assignment,
        and lower temporal variance than image-space history on camera motion.
  \item \textbf{M2 DLSS-class resolve prototype:} $M_s=2$--$8$M,
        $N=3$--$8$M, $K_p=16$--$32$ analytic or $K_p\le 96$ with a small
        neural residual. Demonstrate TAA/SR/frame-interpolation benefit from
        one shared state rather than separate feature-specific histories.
\end{enumerate}

A negative result is useful if it is dimensional: too much $M$ bandwidth, too
large an $N$ retention factor, or too wide a $K$ for the quality gained. The
experiments should report all three numbers for every image-quality result.

\subsection{Mixed-Precision Gaussian State and FP16/FP8/FP4 Policy}

vkGSplat should use lower precision aggressively, but only with a two-copy
rule: keep the canonical estimator state accurate enough to remain stable, and
materialize lower-precision tile packets for bandwidth- and tensor-core-heavy
work. This is different from quantizing the whole renderer. Geometry, motion,
visibility, and estimator control should not become FP4-only state.

The hardware reason is strong. Hopper-class Tensor Cores expose FP8 paths with
E4M3 and E5M2 formats, where E4M3 gives more precision and E5M2 gives wider
range \cite{nvidia2022hopper,nvidia2026transformerenginefp8}. CUDA also exposes
FP4, FP6, FP8, FP16, and BF16 math/intrinsic support in the low-precision API
surface \cite{nvidia2026cudamath}. Blackwell-style NVFP4 uses an E2M1 element
with hierarchical block scaling, making it a plausible tensor format for
neural weights, activations, and compact learned features, but not a safe
standalone geometry format \cite{nvidia2026nvfp4}. The OCP MX specification is
also standardizing microscaled FP4/FP6/FP8-style interchange formats, which is
relevant if vkGSplat later targets non-CUDA accelerators \cite{ocp2026mx}.

The precision policy should be field-specific:

\begin{center}
\small
\begin{tabular}{p{0.30\linewidth}p{0.17\linewidth}p{0.40\linewidth}}
\toprule
State or packet field & Preferred format & Rationale \\
\midrule
Camera matrices, canonical centers $\mu$, covariance update accumulators & FP32 & Keeps reprojection, assignment, and long-lived estimator state stable. \\
Canonical covariance factors, normal, velocity, opacity, radiance moments & FP32/FP16 & Store master accumulators in FP32 where drift matters; expose FP16 render views. \\
Per-tile projected features, analytic resolver features, radiance-cache feature vectors & FP16/FP8 & Bandwidth-heavy and tensor-friendly; use per-tile or per-block scaling. \\
Neural resolver weights, latent residual channels, compressed cache codes & FP8/FP4 & Best FP4 target; require calibration or quantization-aware training. \\
Primitive IDs, material IDs, tile indices, flags, sample counts, timestamps & integer & Never float-quantize identity or ordering metadata. \\
Variance, confidence, visibility gates, disocclusion scores & FP16/FP32 & These control history acceptance; quantization error becomes ghosting. \\
\bottomrule
\end{tabular}
\end{center}

The implementation should therefore split persistent state from working state:

\[
  S^{master}_j = (\mu_j,\Sigma_j,\bar L_j,\bar n_j,\bar v_j,\nu_j,meta_j)
  \quad\rightarrow\quad
  P^{tile}_{j,t} = Q_{p,t}(S^{master}_j),
\]

The master copy uses FP32/FP16 plus integer metadata. The tile packet is a
transient low-precision view used by projection, feature resolve,
radiance-cache lookup, or neural decode. For a block $b$ of feature values $x_i$, a generic
block-scaled packet is

\[
  q_i = Q_p(x_i/s_b),
  \qquad
  \tilde{x}_i = s_b D_p(q_i),
\]

where $p\in\{\mathrm{fp8},\mathrm{fp4}\}$, $Q_p$ quantizes, $D_p$ dequantizes,
and $s_b$ is a tile, cache-page, or micro-block scale. NVFP4-style packets add
a tensor-level scale:

\[
  \tilde{x}_i = s_{global}s_{block}D_{E2M1}(q_i).
\]

This is a natural match for Gaussian radiance cache pages: the page table can
store both residency metadata and quantization scales, just as KV-cache systems
store block/page metadata for attention. The page size should be chosen so one
cache page holds a small integer number of Gaussian feature blocks and their
scales, for example $16$--$64$ Gaussians times $K_h=16$--$64$ feature lanes.

Concrete CUDA kernels:

\begin{enumerate}
  \item \textbf{PackGaussianTilePackets:} read canonical SoA state, project or
        select tile-local Gaussians, compute per-block scales, and emit FP16,
        FP8, or FP4 packet tensors plus scales.
  \item \textbf{TensorFeatureProjectFP8:} evaluate small feature projections or
        MLP-style basis transforms with FP8 inputs and FP16/FP32 accumulation.
  \item \textbf{QuantizedRadianceCacheLookup:} gather page-resident radiance
        features, dequantize into registers/shared memory, and blend with
        analytic visibility and confidence gates.
  \item \textbf{NeuralResolveFP4FP8:} run the optional transformer/MLP resolver
        on FP8 activations and FP4/FP8 weights, accumulating in FP16 or FP32.
  \item \textbf{ApplyGaussianDeltas:} write updates back to the canonical
        state in FP32/FP16 only after compatibility gates and reductions.
\end{enumerate}

The first experiments should treat precision as an ablation axis rather than a
one-way optimization:

\begin{center}
\small
\begin{tabular}{p{0.20\linewidth}p{0.50\linewidth}p{0.20\linewidth}}
\toprule
Ablation & Change & Pass condition \\
\midrule
P0 & FP32 canonical state and FP32/FP16 resolve baseline & Reference. \\
P1 & FP16 render view for color, opacity, normal, velocity, and covariance factors & $\le 0.2$dB PSNR loss and no visible increase in flicker. \\
P2 & FP8 projected features and radiance-cache feature pages & $\le 0.3$dB loss, $\le 10\%$ flicker increase, measurable bandwidth reduction. \\
P3 & FP4 latent resolver weights or residual channels only & $\le 0.5$dB loss versus FP8 resolver and at least $1.3\times$ resolver throughput or bandwidth gain. \\
P4 & FP4 geometry/covariance negative control & Expected to fail; confirms why canonical geometry remains higher precision. \\
\bottomrule
\end{tabular}
\end{center}

Report image metrics (PSNR, SSIM, LPIPS), temporal metrics (flicker,
disocclusion, interpolation endpoint error), and GPU metrics (cache bandwidth,
tensor-core utilization, packet conversion cost, and memory footprint). A
successful result would show that FP16/FP8 reduce the cost of the shared
Gaussian state without weakening history rejection. FP4 should only become
useful once the work has been turned into a calibrated tensor decode problem.

\subsection{Transformer Architecture: Optional Resolver, Not Core State}

Transformer architecture is relevant, but it should enter vkGSplat at the
resolver boundary rather than at the core Gaussian-state boundary. NVIDIA's
DLSS 4 moved Super Resolution, Ray Reconstruction, and DLAA from a CNN lineage
to a real-time vision-transformer model with self-attention over frame and
multi-frame evidence \cite{nvidia2025dlss4}. NVIDIA's DLSS 4.5 announcement
then describes a second-generation transformer model for Super Resolution
\cite{nvidia2026dlss45,nvidia2026dlssdeveloper}. This is strong evidence that
a DLSS-class product should expect transformer-style neural resolving to be
competitive.

For vkGSplat, however, a transformer should not replace the explicit Gaussian
state. The project hypothesis is that geometry, visibility, radiance, motion,
and confidence should be maintained explicitly in CUDA tensors. A transformer,
if used, consumes the projected Gaussian feature tensors and predicts a bounded
residual:

\[
  \hat I_{hi} = R_\theta(F_t, F_{t-1}, D_t, V_t, M_t) + I_{analytic},
\]

where $F_t\in\mathbb{R}^{B\times H_o\times W_o\times K_p}$ is the projected
Gaussian feature tensor, $D_t$ is depth, $V_t$ is motion, $M_t$ is a mask/ID
feature tensor, and $I_{analytic}$ is the non-neural Gaussian/TAA/SR resolve.
The model should learn what the analytic state cannot model cleanly:
high-frequency shading residuals, difficult disocclusions, specular flicker,
and perceptual anti-aliasing. It should not be responsible for discovering
basic geometry or motion from scratch.

\textbf{Operational placement.} The transformer is used only in the decode
resolver, after the Gaussian state has already been updated and projected. It
is not part of the prefill renderer, not part of Gaussian assignment/update, and
not part of the visibility or radiance-cache acceptance gates. Those stages
must remain inspectable CUDA kernels so the no-transformer baseline is always
valid.

\begin{center}
\small
\begin{tabular}{p{0.25\linewidth}p{0.18\linewidth}p{0.45\linewidth}}
\toprule
Pipeline stage & Transformer? & Reason \\
\midrule
Prefill/acquisition & No & Produce factual seed buffers: color, depth, IDs, normals, motion, and variance. \\
Gaussian update/cache & No & Maintain geometry, visibility, radiance, confidence, and lifetime as explicit state. \\
Feature projection & No model & Project Gaussian and engine features into tile/window tensors for the resolver. \\
Decode resolver & Optional yes & Predict bounded residuals, confidence repair, SR detail, or disocclusion fill. \\
Final composite & No & Apply residuals and masks to the analytic resolve with clamps and debug counters. \\
\bottomrule
\end{tabular}
\end{center}

In kernel terms, the optional path is:

\[
\begin{gathered}
\texttt{ProjectFeatures}
\rightarrow \texttt{PackResolverTokens}
\rightarrow \texttt{RunTransformerResolver}\\
\rightarrow \texttt{CompositeResolverOutput}.
\end{gathered}
\]

The token tensor can be written as

\[
  Z_t = \mathrm{PackTokens}(F_t,A_t,H_{t-k:t-1}),
\]

where $A_t$ is the current seed bundle and $H_{t-k:t-1}$ are motion-aligned
history features. The resolver predicts bounded corrections:

\[
  (\Delta I_t,\Delta c_t,\Delta m_t)=T_\theta(Z_t),
  \qquad
  \hat I_t = I_{analytic} + \mathrm{clamp}(\Delta I_t).
\]

Here $\Delta c_t$ can adjust confidence and $\Delta m_t$ can repair masks or
holes. The clamps are not cosmetic: they enforce that the transformer is a
bounded decoder over explicit state, not the owner of geometry or temporal
truth. Disabling the transformer should reduce to the analytic CUDA resolve
without changing the prefill, Gaussian state, cache layout, or metrics.

A practical transformer budget should be windowed and temporal-local rather
than global full-frame attention. Full attention over $P$ pixels scales as
$O(P^2)$ and is not a credible real-time starting point. Window attention with
window size $w\times w$ scales as

\[
  O(P\,w^2\,K_p),
\]

plus small cross-frame token mixing. For 1080p, $w=8$ gives 64-token local
attention; $w=16$ gives 256 tokens and should be treated as a high-quality
ablation, not the default. The first neural resolver should therefore be:

\begin{itemize}
  \item a tiny windowed transformer or hybrid Conv/Transformer block,
  \item input $K_p=32$--$96$ projected Gaussian and engine features,
  \item two to four previous-frame feature tensors at most,
  \item output either residual color, confidence/mask correction, or SR color,
  \item trained and evaluated against an equal-parameter image-space model.
\end{itemize}

The transformer decision should be staged:

\begin{enumerate}
  \item \textbf{No-transformer baseline:} analytic CUDA tensors only. This must
        already demonstrate that shared Gaussian state improves at least one of
        TAA, SR, denoising, radiance caching, or interpolation.
  \item \textbf{Small transformer resolver:} consume projected Gaussian feature
        tensors and compare against an image-only transformer with the same
        parameter count and input resolution.
  \item \textbf{Product-scale transformer:} only after the shared state wins the
        ablation. This path may resemble DLSS-class neural resolving, but it is
        not needed to validate the vkGSplat integration hypothesis.
\end{enumerate}

The pass/fail rule is simple. If a transformer improves both image-space and
Gaussian-feature versions equally, the value is mostly neural resolving. If the
same transformer is smaller, faster, or more stable when fed projected Gaussian
features, then the explicit Gaussian state is helping. Record $M$, $N$, $K_p$,
parameter count, attention window size, and frame history length for every
neural experiment.

\subsection{Why Transformer/LLM Architecture Reuse Is Plausible}

The striking architectural point is that vkGSplat can look like graphics at the
API boundary while looking like Transformer serving inside the expensive part of
the machine. The prefill stage acquires structured evidence from the scene:
color, depth, IDs, barycentrics, normals, motion, visibility, and sparse
radiance samples. The decode stage repeatedly queries persistent state, pages
Gaussian/radiance features, projects local tokens, and resolves pixels, higher
resolution samples, or intermediate frames. That is much closer to a
stateful cache-plus-tensor workload than to a classical fixed-function graphics
pipeline.

With no learned model, the analogy is still useful because Gaussian history and
radiance cache pages behave like a structured memory. With an optional
transformer resolver, the analogy becomes literal: projected Gaussian features,
tile features, and motion-aligned history become tokens, while the Gaussian
radiance cache plays the role of a geometry-aware KV cache. This does not mean
vkGSplat is an LLM. It means the hardware primitives that make LLM serving fast
-- tiled matrix engines, low-precision formats, cache paging, block tables,
gather/scatter, reductions, and bandwidth scheduling -- are also the right
primitives for the decode half of this renderer.

\begin{center}
\small
\begin{tabular}{p{0.23\linewidth}p{0.32\linewidth}p{0.34\linewidth}}
\toprule
Transformer/LLM serving idea & vkGSplat counterpart & Hardware implication \\
\midrule
Prompt prefill & Raster/ray/splat acquisition produces seed evidence $A_t$ & Keep graphics-specific logic narrow and evidence-oriented. \\
KV cache & Persistent Gaussian history, radiance cache, motion state, and feature pages & Invest in SRAM/HBM bandwidth, paging, and cache residency policy. \\
Token blocks/pages & Tile pages, Gaussian pages, radiance-cache pages, and feature blocks & Reuse block tables, page allocators, and compaction/eviction logic. \\
Attention over tokens & Local attention or analytic mixing over pixels, Gaussians, tiles, and history frames & Favor block-sparse tensor kernels and tile-local reuse. \\
Decode loop & Per-frame TAA/SR/denoise/framegen resolve from persistent state & Optimize repeated cache queries and low-launch-count decode kernels. \\
Quantized KV cache & FP8/FP4 Gaussian feature pages and neural resolver packets & Reuse low-precision scaling, calibration, and packet formats. \\
Continuous batching & Batched cameras, views, scenes, or training-data jobs & Schedule many small tile/view jobs like serving requests. \\
Speculative decoding & Motion-predicted intermediate frames or candidate disocclusion fills & Add cheap prediction paths with later validation/rejection. \\
\bottomrule
\end{tabular}
\end{center}

A useful mental diagram is:

\begin{figure}[h]
\centering
\small
\setlength{\fboxsep}{6pt}
\begin{tabular}{c}
\fbox{\begin{minipage}{0.78\linewidth}
\centering
\textbf{Graphics/evidence prefill}\\
minimal raster, ray, or splat acquisition\\
$A_t=(L,D,ID,X,N,V,M,\nu,\sigma^2)$
\end{minipage}}\\[0.6em]
$\Downarrow$\\[0.6em]
\fbox{\begin{minipage}{0.78\linewidth}
\centering
\textbf{Persistent Gaussian/cache state}\\
$G_t$ plus paged radiance/history/features $\mathcal{C}_t$\\
page tables, lifetimes, confidence, scales, visibility gates
\end{minipage}}\\[0.6em]
$\Downarrow$\\[0.6em]
\fbox{\begin{minipage}{0.78\linewidth}
\centering
\textbf{Transformer-like decode}\\
analytic or learned tile resolver\\
TAA, SR, denoise, radiance reuse, frame interpolation, frame generation
\end{minipage}}
\end{tabular}
\caption{vkGSplat as structured evidence acquisition plus stateful decode. The
prefill side benefits from a small amount of graphics-specific visibility
hardware. The decode side is dominated by tensor, cache, paging, quantization,
and scheduling machinery that overlaps strongly with Transformer serving.}
\label{fig:vkgsplat-llm-serving-analogy}
\end{figure}

This leads to a concrete hardware rule. If area is scarce, graphics-specific
hardware should be added only where it improves evidence quality per joule:
depth, primitive IDs, barycentrics, motion vectors, simple texture/material
samples, and perhaps ray-hit evidence. The larger investment should remain in
the decode substrate: tensor engines, SRAM/HBM bandwidth, page tables,
low-precision feature packets, tile scheduling, and sparse gather/reduce
kernels. A full legacy raster/ROP/texture/RT stack is less attractive unless
experiments show that the prefill evidence is the dominant quality or energy
bottleneck.

The falsifiable experiment is an area-proxy ablation, not just an image-quality
ablation. Hold total time or bandwidth budget fixed and compare:

\begin{enumerate}
  \item \textbf{graphics-heavy prefill, weak decode:} higher-spp or higher-res
        seed buffers, but small Gaussian/cache decode;
  \item \textbf{minimal prefill, strong decode:} sparse seed buffers, larger
        Gaussian state, better cache paging, and tensorized decode;
  \item \textbf{minimal prefill plus transformer resolver:} same seed and state
        as (2), but add a small learned tile resolver;
  \item \textbf{balanced bridge:} minimal prefill plus a narrow tile-splat
        accelerator for Gaussian binning, projection, and compositing.
\end{enumerate}

If (2) or (3) wins at equal bandwidth and latency, vkGSplat should lean toward
LLM-style tensor/cache architecture. If (1) wins decisively, then more graphics
prefill hardware is justified. If (4) wins, the right hardware is not a full
classic graphics pipeline but a narrow Gaussian/tile bridge between acquisition
and decode.

\subsection{Prefill/Decode Split and Hardware Investment}

vkGSplat should explicitly adopt a prefill/decode architecture. The prefill
stage is the expensive real-frame evidence stage; the decode stage is the
cheaper reconstruction/generation stage that reuses persistent Gaussian and
feature caches.

\[
\begin{aligned}
  \textsc{Prefill}(A_t,G_{t-1}) &\rightarrow (G_t,\mathcal{C}_t),\\
  \textsc{Decode}(G_t,\mathcal{C}_t,\alpha,r) &\rightarrow \hat I_{t+\alpha}^{(r)}.
\end{aligned}
\]

Here $A_t$ is the acquired real-frame evidence tensor bundle, $G_t$ is the
persistent Gaussian state, $\mathcal{C}_t$ is the tile/feature/cache bundle,
$\alpha$ is the generated-frame time offset, and $r$ is the super-resolution
scale. The prefill stage performs acquisition, hit extraction, assignment,
Gaussian update, visibility decay, grid build, tile-list construction, and
optional neural key/value construction. The decode stage projects/querys this
state to produce TAA, SR, DLAA, frame interpolation, or generated frames.

The average cost of one real frame plus $q$ decoded outputs is approximately

\[
  T_{avg}(q)=\frac{T_{prefill}+qT_{decode}}{q+1}.
\]

The architecture is useful only if $T_{decode}\ll T_{prefill}$ and cache
invalidations remain spatially sparse. A rough cost model is

\[
\begin{aligned}
  T_{prefill} &\sim aM_{raw}+bM_sC+cN+T_{scan/sort}+T_{state},\\
  T_{decode}  &\sim dP_{out}K_p+eN_{vis}+T_{attn/cache},\\
  T_{attn/cache} &\sim fP_{out}w^2K_a
        \quad\text{for windowed attention.}
\end{aligned}
\]

This model clarifies the hardware bet. Rasterization, texture, and ray-tracing
hardware remain essential evidence engines. They generate primary visibility,
texture/material samples, ray-traced lighting evidence, shadows, reflections,
disocclusion masks, and correctness anchors. Tensor hardware and KV-cache-like
memory systems are the likely decode engines: they turn sparse explicit state
and projected feature tensors into high-rate frames, super-resolution outputs,
and neural residuals.

Therefore vkGSplat should not frame the choice as raster/texture/RT
\emph{versus} tensor/KV cache. The better architecture is:

\begin{quote}
Rasterization, texture, and ray tracing produce high-quality prefill evidence;
tensor cores, cache residency, sparse lookup, and KV-like feature caches make
repeated decode cheap.
\end{quote}

For vkGSplat research effort, the marginal investment should lean toward the
second half because ordinary rasterization already exists while DLSS-class
multi-decode reconstruction depends on cache reuse. A practical effort split is:

\begin{center}
\begin{tabular}{lll}
\toprule
Area & Suggested emphasis & Reason \\
\midrule
Raster/texture/RT acquisition & 20--30\% & reliable evidence and baselines \\
Tensorized Gaussian/cache kernels & 40--50\% & core vkGSplat differentiator \\
KV-like neural/tile feature cache & 15--25\% & makes multi-decode economical \\
Interop/profiling/evaluation & 10--15\% & proves real frame-time value \\
\bottomrule
\end{tabular}
\end{center}

For hardware architecture, the recommendation is similar but more balanced:
keep improving raster, texture, and ray traversal enough to feed accurate
prefill frames, but put the speculative DLSS-5-class upside into tensor
throughput, low-latency cache residency, sparse/tiled feature lookup, fast
quantized feature storage, and synchronization between graphics and compute.
A Gaussian-aware cache is not literally an LLM KV cache, but it plays an
analogous role: expensive evidence is encoded once, then many decode queries
reuse compact keys, values, confidence, visibility, and motion features.

A concrete cache layout for experiments is:

\[
  \mathcal{C}_t =
  (\texttt{tile\_ranges},\texttt{gaussian\_indices},K_t,V_t,M_t),
\]

where $K_t\in\mathbb{R}^{T\times B_t\times K_a}$ stores per-tile keys,
$V_t\in\mathbb{R}^{T\times B_t\times K_v}$ stores projected Gaussian/radiance
values, $M_t$ stores masks, depth ranges, IDs, and confidence, $T$ is the
number of tiles, and $B_t$ is a capped per-tile cache budget. Decode then
queries only the local tile and neighbor tiles, not the entire Gaussian set.

The pass/fail experiments should be:

\begin{enumerate}
  \item \textbf{Prefill amortization:} compare monolithic update+resolve against
        prefill once plus 1, 2, 4, and 6 decodes. Report $T_{avg}(q)$,
        cache memory, ghosting, and disocclusion error.
  \item \textbf{Cache quantization:} compare fp32, fp16, fp8/int8-like cached
        features for $K_t,V_t$ at fixed quality.
  \item \textbf{KV-cache ablation:} decode from raw projected Gaussians, from
        tile feature caches, and from a small transformer with cached keys and
        values.
  \item \textbf{Partial refresh:} invalidate only tiles/primitives with large
        disocclusion or motion residuals and measure whether quality recovers
        without full prefill.
\end{enumerate}

The decision rule is: invest more in tensor/KV-cache optimization if prefill
amortization gives at least $2\times$--$4\times$ displayed-frame throughput at
acceptable temporal error, and if cache memory stays within the $N,K$ budgets
above. If decode quality collapses without frequent full prefill, invest more
in acquisition and motion/visibility hardware instead.

\subsection{Reusing KV-Cache Optimization for Gaussian Radiance Cache}

KV-cache hardware and software optimization is directly relevant to vkGSplat,
but the reuse is architectural rather than literal. LLM inference caches keys
and values for previously processed tokens so repeated decode steps avoid
recomputing the prefix. A Gaussian radiance cache should cache spatial,
visibility-aware keys and radiance/history values so repeated TAA/SR/framegen
decodes avoid re-querying or re-updating the entire Gaussian state.

\begin{center}
\small
\begin{tabular}{p{0.22\linewidth}p{0.34\linewidth}p{0.30\linewidth}}
\toprule
LLM KV cache & Gaussian radiance cache & Reusable idea \\
\midrule
Token position & tile/cell/Gaussian address & logical-to-physical block map \\
Key vector & position, normal, depth, ID, motion, confidence & query-key scoring \\
Value vector & radiance, SH/lobe moments, variance, age & cached decode payload \\
Prefix reuse & static or slow scene regions & sharing across generated frames \\
Paged blocks & tile/Gaussian pages & low-fragmentation dynamic allocation \\
KV quantization & radiance/feature quantization & bandwidth and residency reduction \\
Copy-on-write & page update on visibility/material change & sparse refresh \\
\bottomrule
\end{tabular}
\end{center}

PagedAttention shows that dynamic KV caches benefit from block paging, reduced
fragmentation, and flexible sharing across requests \cite{kwon2023pagedattention}.
vAttention argues for retaining contiguous virtual memory while allocating
physical memory on demand, avoiding kernel complexity caused by explicitly
paged layouts \cite{prabhu2025vattention}. vkGSplat should test both ideas in
Gaussian form: explicit block tables are easier to prototype and inspect, while
contiguous virtual addressing may be better if CUDA virtual-memory APIs keep
kernels simpler.

A Gaussian radiance page should have a logical address

\[
  a = (tile, depth\_layer, primitive\_group, material\_group, epoch),
\]

and a physical page containing a capped number $B_p$ of cache records:

\[
  \mathcal{P}_a =
  (K_a,V_a,M_a,E_a),
\]

where

\[
\begin{aligned}
  K_a &\in \mathbb{R}^{B_p\times K_{key}},\\
  V_a &\in \mathbb{R}^{B_p\times K_{val}},\\
  M_a &\in \mathbb{Z}^{B_p\times K_{meta}},\\
  E_a &\in \mathbb{R}^{B_p\times K_{err}}.
\end{aligned}
\]

The key tensor should include compact geometry and identity features:
position relative to tile, depth range, normal cone, primitive/material ID
hashes, velocity, confidence, and visibility age. The value tensor should store
radiance and reconstruction payloads: RGB or SH coefficients, indirect-radiance
moments, lobe/roughness bins, variance, temporal confidence, and optional
neural features. The error tensor stores residual statistics used for eviction,
refresh, and quantization health.

The decode query becomes local attention or local compatibility over a small
set of pages:

\[
  \hat L(q)=
  \sum_{a\in\mathcal{N}(tile(q))}\sum_{j\in\mathcal{P}_a}
  \mathrm{softmax}_j\left(\frac{\phi(q)^\top K_{a,j}}{\sqrt{K_{key}}}
  -\lambda_d\Delta d-\lambda_n\Delta n-\lambda_{id}\Delta id\right)V_{a,j}.
\]

The hard visibility and identity gates remain outside the softmax. This is a
major difference from LLM KV cache: wrong reuse in text usually changes token
probability smoothly; wrong reuse in rendering creates obvious ghosting,
light leaks, and depth-edge errors. Therefore the page table must be
geometry-aware:

\begin{itemize}
  \item \textbf{Hard invalidation:} primitive/material mismatch, large depth
        residual, disocclusion, or topology epoch mismatch invalidates records.
  \item \textbf{Soft decay:} occlusion age, variance, and motion residual lower
        confidence without immediately freeing the page.
  \item \textbf{Copy-on-write:} static shared pages can serve multiple decode
        outputs until a prefill update touches their primitive/tile group.
  \item \textbf{Motion prefetch:} tile pages are prefetched along screen-space
        and Gaussian velocities for frame interpolation.
\end{itemize}

Quantization should be asymmetric. Geometry, depth bounds, primitive identity,
and confidence should stay fp16/fp32 or integer-exact until experiments prove
otherwise. Radiance values, SH/lobe features, and neural features are better
candidates for fp16, fp8, int8, or product/vector quantization. The right first
experiment is not maximum compression; it is measuring when quantized cached
radiance begins to cause visible flicker, color bias, or light leakage:

\[
  Q_{safe} = \min_Q
  \{\mathrm{bytes}(Q):\Delta E_{temporal}<\epsilon_t,
                    \Delta E_{edge}<\epsilon_e,
                    \Delta E_{radiance}<\epsilon_L\}.
\]

The CUDA kernels implied by this design are:

\begin{enumerate}
  \item \texttt{BuildRadiancePageKeys}: write page keys from active Gaussians,
        tile lists, depth ranges, and primitive/material groups.
  \item \texttt{AllocateRadiancePages}: allocate physical pages from a device
        pool using a block table or virtual-memory-backed allocator.
  \item \texttt{ScatterRadianceValues}: write radiance, variance, confidence,
        and optional neural values into page-local SoA tensors.
  \item \texttt{DecodeRadiancePages}: query local pages during TAA/SR/framegen
        resolve using hard gates plus local key-value scoring.
  \item \texttt{InvalidateRadiancePages}: clear or decay pages after
        disocclusion, primitive change, or large residual error.
  \item \texttt{QuantizeRadiancePages}: pack stable value pages to fp8/int8-like
        formats and leave unstable pages at higher precision.
\end{enumerate}

This section turns the hardware question into a concrete vkGSplat thesis:

\begin{quote}
The reusable hardware idea is not ``attention for everything.'' It is paged,
quantized, high-residency, queryable state for repeated decode. Gaussian
radiance cache is the graphics analogue of KV cache, with geometry-aware
invalidation added.
\end{quote}

The pass/fail experiments are:

\begin{enumerate}
  \item \textbf{Paged vs contiguous cache:} compare explicit page tables,
        contiguous virtual allocation, and simple flat arrays. Report page
        waste, allocator time, decode time, and kernel complexity.
  \item \textbf{Radiance-cache quantization:} fp32/fp16/fp8/int8-like values at
        fixed $N,K$ and fixed frame sequence. Report temporal flicker,
        radiance bias, and edge leakage.
  \item \textbf{KV-style sharing:} reuse static pages across 1, 2, 4, and 6
        decode outputs. Report copy-on-write rate and stale-page error.
  \item \textbf{Geometry invalidation:} deliberately introduce disocclusions,
        moving occluders, animated materials, and thin geometry. The cache
        passes only if hard gates prevent light leaks and ghost trails.
\end{enumerate}

The investment decision is positive if block-paged or virtual-memory-backed
radiance cache reduces memory waste and decode bandwidth without increasing
visible stale-history artifacts. It is negative if geometry-aware invalidation
forces so many page refreshes that the system behaves like full prefill every
frame.

\subsection{CUDA Data Layout and Binning}

The first reconstruction implementation should be CUDA-first, but it should
not depend on a CUDA acquisition backend existing yet. Recorded seed frames and
CPU/Metal-produced seed frames are valid inputs as long as they satisfy the
same device view. Once CUDA acquisition exists, it can alias or write those
buffers directly without changing the reconstruction kernels.

Avoid a global nearest-neighbor search. Use a GPU-built hashed grid keyed by
quantized position plus primitive ID:

\[
  h(x,id) =
  (\lfloor x_x / h_c \rfloor,
   \lfloor x_y / h_c \rfloor,
   \lfloor x_z / h_c \rfloor,
   id).
\]

For each sample, search the cell and its 26 neighbors. Cap candidates per cell:

\[
  |B_h| \le B_{max}.
\]

If a cell overflows, reject low-confidence Gaussians first. If overflow
persists, split the cell by primitive ID or material ID. The data layout should
be structure-of-arrays so projection and update kernels can coalesce loads:

\[
  G =
  (\mu_x,\mu_y,\mu_z,\Sigma_0,\ldots,\Sigma_5,
   c_x,c_y,c_z,\nu,v,m,id_{prim},id_{mat},cell).
\]

The CUDA staging buffers are:

\begin{itemize}
  \item \textbf{sample buffers:} one entry per seed pixel/path event,
  \item \textbf{cell counts:} atomically incremented during binning,
  \item \textbf{cell offsets:} prefix-summed with CUB or an equivalent scan,
  \item \textbf{cell members:} compact Gaussian indices sorted or grouped by
        cell,
  \item \textbf{tile lists:} projected Gaussian ranges for TAA/SR/framegen,
  \item \textbf{debug counters:} gate rejections, atomics, overflows, clamps,
        prunes, and stale-history energy.
\end{itemize}

The prototype should measure atomics and memory traffic from the beginning.
The research question is not only whether the equations work; it is whether
the shared state maps naturally to CUDA.

\subsection{CUDA Hash Grid Build}

Build the hash grid with a standard count/scan/scatter pipeline:

\begin{enumerate}
  \item \texttt{ComputeCellKeys}: one thread per Gaussian computes
        $h(\mu_j,id_j)$ and writes a compact cell key.
  \item \texttt{CountCells}: atomically increments a cell count.
  \item \texttt{ExclusiveScan}: converts counts to cell offsets.
  \item \texttt{ScatterCellMembers}: writes Gaussian indices into
        \texttt{cell\_members}.
  \item \texttt{ClampOverflow}: marks excess low-confidence members for prune,
        split, or delayed update.
\end{enumerate}

The per-cell occupancy distribution is a first-order performance signal:

\[
  \bar{B}=\frac{1}{N_c}\sum_h |B_h|,
  \qquad
  B_{p99}=\mathrm{percentile}_{99}(|B_h|).
\]

The method is healthy only if $B_{p99}$ stays bounded under camera motion and
moderate scene complexity. If a few cells dominate, assignment becomes
divergent and tile projection develops long-tail latency.

\subsection{CUDA Assignment and Update Strategy}

Avoid updating Gaussian state directly from sample threads when possible. The
simple atomic path is useful for M0, but the intended CUDA path is two-stage:
sample-parallel assignment followed by Gaussian-parallel reduction.

Sample assignment writes:

\[
  A_i=(j_i^*, W_{ij_i^*}, x_i,c_i,L_i,n_i,u_i,\rho_i).
\]

Then reduce by Gaussian ID:

\[
  S_j = \sum_{i:j_i^*=j} W_{ij},
  \qquad
  X_j = \sum_{i:j_i^*=j} W_{ij}x_i,
  \qquad
  C_j = \sum_{i:j_i^*=j} W_{ij}c_i.
\]

The update kernel becomes one thread block or warp group per active Gaussian:

\[
  \mu_j' =
  \frac{m_j\mu_j+X_j}{m_j+S_j+\epsilon},
  \qquad
  c_j' =
  \frac{m_jc_j+C_j}{m_j+S_j+\epsilon}.
\]

Radiance, velocity, and variance use the same reduced sums. This avoids many
floating-point atomics to the same Gaussian and makes update cost depend on
active Gaussians rather than raw samples. The first implementation can keep a
compile-time switch:

\begin{itemize}
  \item \textbf{atomic update path:} fastest to implement, useful for small
        scenes and debugging,
  \item \textbf{reduced update path:} intended performance path, using sorted
        assignments or bucketed per-cell reductions.
\end{itemize}

The experiment should report both. If the reduced path is required for quality
because atomics create unstable history, that is acceptable. If it is required
but too expensive, the shared-state CUDA design needs a different addressing
scheme.

\subsection{CUDA Spawn, Prune, Split, and Merge}

Spawning should be queue-based:

\[
  q_i =
  \begin{cases}
  1, & W_{ij^*}<\epsilon_W,\\
  0, & \text{otherwise}.
  \end{cases}
\]

Prefix-sum $q_i$ to allocate new Gaussian slots from a device free list:

\[
  j_{new}=F[\mathrm{atomicAdd}(F_{head},1)].
\]

Pruning writes dead Gaussian indices back to the free list:

\[
  dead_j =
  [\nu_j < \epsilon_\nu]\lor[t-\tau_j>T_{max}].
\]

Split and merge should be separate compaction passes, not mixed into the
sample update kernel. Split high-variance Gaussians by emitting two children
along the largest tangent eigenvector:

\[
  \mu_{j\pm}=\mu_j\pm \kappa\sqrt{\lambda_{max}}e_{max},
  \qquad
  \Sigma_{j\pm}=\eta_{split}\Sigma_j.
\]

Merge is a low-frequency maintenance pass over local cells. The CUDA kill
condition is simple: if split/merge must run every frame to avoid blur or
explosive memory growth, the state representation is too fragile.

\subsection{Tile List Construction and Reuse}

Projection should use one tile list for all image-space consumers. For each
Gaussian, compute its projected ellipse and conservative bounding box:

\[
  B_j =
  [p_j-r_j, p_j+r_j],
  \qquad
  r_j = k_\sigma
  \sqrt{\lambda_{max}(\Sigma_{screen,j})}.
\]

The tile coverage is:

\[
  T_j =
  \{\,\tau : \tau \cap B_j \neq \emptyset\,\}.
\]

Build tile lists with count/scan/scatter, mirroring the hash-grid build:

\[
  count_\tau = \sum_j [\tau \in T_j].
\]

Then TAA, SR, feature projection, and frame generation consume the same
\texttt{tile\_members}. The consumers differ only in output resolution,
confidence gates, and payload channels. A block-per-tile projection kernel can
load a chunk of Gaussian indices into shared memory, evaluate per-pixel
weights, and accumulate:

\[
  O_p =
  \frac{\sum_{j\in tile(p)} a_{pj} y_j}
       {\sum_{j\in tile(p)} a_{pj}+\epsilon}.
\]

The crucial CUDA measurement is tile reuse:

\[
  R_{tile}=
  \frac{N_{consumers}\,N_{tile\_members}}
       {N_{tile\_members}+N_{extra\_members}}.
\]

If $R_{tile}$ is close to one, the shared state is not buying much. If it is
large and quality remains stable, vkGSplat has a concrete integration win.

\subsection{Persistent State Versioning}

Frame interpolation and temporal reconstruction need more than one state view.
Maintain a ring buffer for compact motion-critical fields:

\[
  G_j^{hist}(r) =
  (\mu_j^r,\Sigma_j^r,c_j^r,\nu_j^r,\tau_j^r),
  \qquad
  r=t\bmod R.
\]

The full radiance/path payload can stay single-buffered unless a path-space
experiment needs old values. The CUDA memory rule is:

\[
  M_{history} =
  R\,M_{motion\_core} + M_{payload}.
\]

This is cheaper than duplicating the entire Gaussian state for TAA, SR, and
frame generation. It also gives frame generation a clean source for
$G_j(t_0)$ and $G_j(t_1)$ without keeping separate motion caches.

\subsection{Default Prototype Parameters}

The first demo should use explicit conservative defaults:

\begin{center}
\begin{tabular}{lll}
\toprule
Parameter & Initial value & Purpose \\
\midrule
$\epsilon_W$ & $0.05$ & spawn threshold \\
$\epsilon_d$ & $10^{-3}$--$10^{-2}$ & depth gate \\
$\gamma_n$ & $8$ & normal sharpness \\
$\eta_{id}$ & $0$ & hard primitive rejection \\
$\lambda_{occ}$ & $0.98$ & slow occlusion decay \\
$\lambda_{dis}$ & $0.2$ & fast stale-history decay \\
$\lambda_{frustum}$ & $0.995$ & preserve offscreen state briefly \\
$m_{split}$ & $32$ & split mass threshold \\
$v_{split}$ & scene-dependent & split variance threshold \\
$B_{max}$ & $8$--$16$ & max Gaussians per hash cell \\
\bottomrule
\end{tabular}
\end{center}

These values are not claims. They are a reproducibility anchor so vkGSplat can
compare changes without silently moving the target.

\subsection{Algorithm Summary}

\begin{verbatim}
for each frame t:
    seed = acquire_seed_frame(t)  // CPU, Metal, CUDA RT/raster, or 3DGS
    upload_or_alias_seed_buffers(seed)

    launch ExtractSamples(seed, sample_buffers)
    launch UpdateVisibility(gaussians, seed.depth, seed.primitive_id)
    launch ComputeCellKeys(gaussians, cell_keys)
    launch CountCells(cell_keys, cell_counts)
    scan cell_counts -> cell_offsets
    launch ScatterCellMembers(cell_keys, cell_offsets, cell_members)

    launch AssignSamples(sample_buffers, cell_members, assignments, spawn_flags)
    scan spawn_flags -> spawn_offsets
    launch SpawnGaussians(spawn_offsets, free_list, gaussians)

    if use_reduced_update:
        sort_or_bucket assignments by gaussian_id
        launch ReduceAssignments(assignments, gaussian_deltas)
        launch ApplyGaussianDeltas(gaussian_deltas, gaussians)
    else:
        launch AtomicUpdateGaussians(assignments, gaussians)

    launch PruneAndCompactIfNeeded(gaussians, free_list)
    launch SplitOrMergeMaintenanceIfScheduled(gaussians)

    launch BuildTileLists(gaussians, camera_t, tile_members)
    launch ProjectTAA(tile_members, gaussians, taa_frame)
    launch ProjectSR(tile_members, gaussians, sr_frame)
    launch ProjectFeatures(tile_members, gaussians, features)

    if path_tracing:
        launch QueryRadiance(gaussians, path_queries)
        launch UpdatePathSubslots(path_events, gaussians)

    if interpolation:
        launch SynthesizeGenerated(tile_members, gaussians, alpha, framegen_out)
\end{verbatim}

\subsection{CUDA Kernel Decomposition}

The CUDA prototype should be decomposed into explicit kernels so each
mathematical claim maps to a launch, a memory access pattern, and a measurable
cost:

\begin{center}
\begin{tabular}{llll}
\toprule
CUDA kernel & Input & Output & Main risk \\
\midrule
ExtractSamples & seed frame & $S_t$ & bandwidth \\
BinGaussians & $G_j$ & cell lists & atomics \\
AssignSamples & $S_t$, cells & assignments & divergence \\
UpdateMoments & assignments & updated $G_j$ & write conflicts \\
UpdateVisibility & depth, IDs & $\ell_j^t,\nu_j$ & false stale kills \\
BuildTileLists & projected $G_j$ & tile ranges & overdraw \\
ProjectTAA & tile lists & history/confidence & depth leaks \\
ProjectSR & tile lists, scale & high-res resolve & bandwidth \\
QueryRadiance & path events & $\hat{L}^k$ & incoherent reads \\
UpdatePathSubslots & assignments & $A_j^{path}$ & noisy atomics \\
SynthesizeFrame & $G_j^\alpha$ & generated frame & holes \\
ProjectFeatures & tile lists & feature planes & channel cost \\
\bottomrule
\end{tabular}
\end{center}

The CUDA kernels should be deterministic for fixed input where possible, but
perfect CPU-like determinism is not the objective. The objective is to expose
where the unified Gaussian state pays for itself on the GPU: fewer duplicated
history buffers, fewer independent passes, reusable tile lists, and shared
confidence/visibility logic.

\subsection{CUDA Scheduling}

The preferred frame schedule is:

\begin{enumerate}
  \item seed frame and auxiliary buffers are produced by any acquisition
        backend that satisfies the seed ABI,
  \item \texttt{ExtractSamples} writes a compact sample stream,
  \item \texttt{UpdateVisibility} marks visible/occluded/stale Gaussians,
  \item \texttt{BinGaussians} and \texttt{AssignSamples} update the persistent
        state,
  \item \texttt{BuildTileLists} constructs one projection structure,
  \item TAA, SR, feature projection, and frame generation consume the same tile
        lists,
  \item radiance queries update path subslots asynchronously when path tracing
        samples are available.
\end{enumerate}

This schedule is the real integration claim. If TAA, SR, radiance, and frame
generation require separate binning/projection structures, then 3DGS remains a
collection of useful tricks rather than vkGSplat's central CUDA abstraction. If
they require a specific acquisition renderer, the software stack has failed at
an even earlier boundary.

\subsection{Task-to-Experiment Map}

\begin{center}
\begin{tabular}{llll}
\toprule
Task & Equation family & First experiment & Failure signal \\
\midrule
TAA & $H_t,\beta_p$ & Exp. 1 & shimmer or ghosting \\
2x/4x SR & $O_t^{SR}$ & Exp. 1 & blur or scale-specific state \\
Denoising & $F_L,V_L$ & Exp. 2 & variance unchanged \\
Radiance cache & $\hat{L}^k$ & Exp. 2 & light leaking \\
Path-space CV & $\hat{I}_{CV}$ & Exp. 2 & bias exceeds variance gain \\
Interpolation & $\mu^\alpha,\nu^\alpha$ & Exp. 3 & silhouette error \\
Frame generation & $H_{\alpha},E_{FG}$ & Exp. 3 & confident holes \\
Neural resolve & $\Phi_t,N_\theta$ & Exp. 4 & features ignored \\
\bottomrule
\end{tabular}
\end{center}

This table is the falsifiability guardrail. If an equation does not land in an
experiment, it should not be in the first implementation.

\section{Feature Integration}

\subsection{Monte Carlo Denoising and Radiance Caching}

Use path samples to update radiance-bearing Gaussian levels:

\begin{itemize}
  \item level 0: directly visible surface evidence,
  \item level 1: first-bounce diffuse or rough indirect radiance,
  \item level 2: deeper low-frequency residual,
  \item residual level: optional control-variate correction.
\end{itemize}

The cache is queried during low-spp rendering and also projected as a denoising
feature. This tests whether radiance history and image history can be one
state.

\subsection{Temporal AA}

The Gaussian state acts as a world-space history buffer. Instead of sampling
previous-frame pixels, the current frame projects persistent Gaussians into the
current camera and accepts them only when depth, primitive ID, normal, velocity,
and confidence agree.

\subsection{Temporal Super-Resolution}

Jittered low-resolution samples update Gaussian footprints. The output image is
resolved at target resolution by splatting projected footprints. The key
question is whether the state preserves subpixel evidence better than a
well-clamped image-space TAAU baseline.

\subsection{Motion-Accurate Frame Interpolation}

For interpolation between $t_0$ and $t_1$, evaluate each moving Gaussian at
$t_\alpha$:

\[
  \mu_j(t_\alpha) =
  (1-\alpha)\mu_j(t_0) + \alpha\mu_j(t_1)
\]

for the first prototype. Later versions use object transforms, splines,
Kalman-style prediction/correction, or flow-guided residuals. Motion is trusted
only when projected forward/backward residuals are small:

\[
  e_j =
  \left\|\Pi_{t_1}(\mu_j(t_1)) -
  \left(\Pi_{t_0}(\mu_j(t_0)) + u_{0\rightarrow1}\right)\right\|.
\]

This is stronger than plausible frame generation: the midpoint should align
with a true rendered middle frame.

\subsection{Hybrid Learned Resolver}

A small neural resolver is optional later work. It can consume both normal
engine buffers and Gaussian features:

\[
  \text{low-res frame} + \text{engine buffers} + \text{Gaussian features}
  \rightarrow
  \text{reconstructed frame}.
\]

This does not replace the explicit state, and it is not required for the core
framework. It tests whether the explicit state reduces the amount of learning
needed.

\section{What Can Approach DLSS 4.5 or DLSS 5}

Likely approachable in vkGSplat:

\begin{itemize}
  \item improved temporal stability over naive reprojection,
  \item 2x temporal SR for controlled scenes,
  \item motion-accurate midpoint interpolation for rigid camera/object motion,
  \item radiance-cache-assisted diffuse indirect reconstruction,
  \item explicit world-space conditioning for a later optional neural
        lighting/material residual,
  \item a hybrid neural resolver that uses Gaussian features, if the
        no-training core already shows signal.
\end{itemize}

Not likely approachable soon:

\begin{itemize}
  \item generic per-game DLSS 4.5 product quality,
  \item DLSS 5 photoreal material/lighting synthesis without a large trained
        model, art controls, and production validation,
  \item robust multi-frame generation for particles, UI, transparency, and
        specular shading without an image-space learned residual,
  \item DLSS-class latency integration without Reflex-like scheduling,
  \item parity with proprietary NVIDIA training and validation infrastructure.
\end{itemize}

The useful vkGSplat question is therefore not ``can 3DGS replace DLSS?'' It is:
can an explicit Gaussian state provide a better conditioning signal than color,
depth, motion vectors, and exposure alone? If yes, vkGSplat becomes a plausible
open research route toward a DLSS-class stack: analytic Gaussian history first,
then, only if useful, a small neural residual that receives projected Gaussian
features.

\section{Metrics and Pass/Fail Rules}

The experiments should not rely on visual inspection alone. The shared-state
hypothesis needs a scorecard that measures quality, stability, integration, and
cost.

\subsection{Image and Temporal Error}

Reference error:

\[
  E_{ref} =
  \frac{1}{|\Omega|}
  \sum_{p\in\Omega}
  \|O_t(p)-R_t(p)\|_2^2.
\]

Temporal instability after reprojection:

\[
  E_{temp} =
  \frac{1}{|\Omega_v|}
  \sum_{p\in\Omega_v}
  \|O_t(p)-O_{t-1}(p+u_{t\rightarrow t-1})\|_1.
\]

Disoccluded pixels are excluded from $\Omega_v$ and scored separately.

\subsection{Ghosting and Disocclusion}

Given a ground-truth or depth/primitive-derived disocclusion mask $D_t(p)$,
ghosting error is:

\[
  E_{ghost} =
  \frac{1}{|D_t|}
  \sum_{p:D_t(p)=1}
  \|O_t(p)-C_t(p)\|_1.
\]

This penalizes stale history in newly revealed regions. A Gaussian method that
improves PSNR but raises $E_{ghost}$ is not acceptable for the first demo.

\subsection{Interpolation Accuracy}

For midpoint interpolation, compare against a true rendered middle frame:

\[
  E_{mid} =
  \frac{1}{|\Omega|}
  \sum_p
  \|O_{t+\frac{1}{2}}^{gen}(p)-R_{t+\frac{1}{2}}(p)\|_2^2.
\]

Silhouette error uses depth or primitive masks:

\[
  E_{sil} =
  1 -
  \frac{|S_{gen}\cap S_{ref}|}{|S_{gen}\cup S_{ref}|}.
\]

The midpoint is motion-accurate only if both $E_{mid}$ and $E_{sil}$ improve
over image blending and screen-space reprojection.

\subsection{Integration Score}

Define a normalized per-feature improvement over the image-space baseline:

\[
  I_f =
  \frac{E_f^{baseline}-E_f^{method}}
       {E_f^{baseline}+\epsilon}.
\]

For features $f\in\{\mathrm{TAA},\mathrm{SR},\mathrm{Denoise},\mathrm{Interp}\}$:

\[
  I_{shared} =
  \frac{1}{4}\sum_f I_f
  - \lambda_C\frac{C_{method}}{C_{baseline}}
  - \lambda_M\frac{M_{method}}{M_{baseline}}.
\]

The shared Gaussian state passes the integration test if:

\begin{itemize}
  \item at least three of four $I_f$ are positive,
  \item no ghosting metric regresses by more than 10\%,
  \item shared-state cost is below separate-Gaussian cost,
  \item shared-state memory is within 1.5x of the image-space baseline for the
        first CUDA demo, or clearly on a path to that bound.
\end{itemize}

\section{Decisive Experiments}

\subsection{Experiment 0: Baseline Reconstruction Stack}

\textbf{Purpose.} Establish honest non-Gaussian baselines.

\textbf{Inputs.} Low-spp or low-resolution seed frames with color, depth, NDC
depth, primitive ID, camera matrices, jitter, and optional motion vectors.

\textbf{Baselines.}

\begin{itemize}
  \item current vkGSplat reprojection,
  \item TAA with neighborhood clamp,
  \item TAAU 2x with jitter,
  \item SVGF-style denoising,
  \item optical-flow interpolation if available.
\end{itemize}

\textbf{Metrics.} PSNR, SSIM, temporal difference, edge shimmer, ghosting after
disocclusion, silhouette error, runtime, and memory.

\textbf{Expected signal.} This defines the bar. Gaussian results are meaningful
only if these baselines are not straw men.

\subsection{Experiment 1: Single Gaussian State, Multiple Outputs}

\textbf{Purpose.} Demonstrate integration directly.

\textbf{Scene.} Static Cornell-style scene with thin geometry, textured
surfaces, camera jitter, and slow camera pan.

\textbf{Method.} Build one Gaussian state from visible samples. Use that same
state to produce:

\begin{itemize}
  \item TAA history,
  \item 2x super-resolution resolve,
  \item variance/confidence feature frame,
  \item disocclusion mask.
\end{itemize}

\textbf{Baselines.}

\begin{itemize}
  \item separate image-space TAA plus TAAU,
  \item Gaussian TAA only,
  \item Gaussian SR only,
  \item shared Gaussian state for both.
\end{itemize}

\textbf{Metrics.}

\begin{itemize}
  \item TAA shimmer reduction,
  \item SR reference error,
  \item ghosting score,
  \item number of duplicated state buffers,
  \item update/query cost.
\end{itemize}

\textbf{Expected success.} The shared Gaussian state performs within 5--10\%
of specialized per-feature Gaussian variants while beating image-space history
on shimmer or disocclusion.

\textbf{Failure case.} If each feature requires incompatible covariance,
confidence, or filtering rules, 3DGS is not a clean integration layer.

\subsection{Experiment 2: Gaussian Radiance Cache Plus TAA}

\textbf{Purpose.} Test whether lighting history and image history can share
the same Gaussian state.

\textbf{Scene.} Diffuse Cornell box plus one rough glossy variant.

\textbf{Method.} Add path sample logging. Store level-1 indirect radiance in
the same Gaussian entries or in path-depth subslots attached to the same entry.
Project color history for TAA and query indirect radiance for low-spp shading.

\textbf{Baselines.}

\begin{itemize}
  \item path tracing plus SVGF,
  \item Gaussian TAA without radiance cache,
  \item separate radiance cache plus separate TAA history,
  \item unified Gaussian state with color and radiance subslots.
\end{itemize}

\textbf{Metrics.} Indirect-light MSE, variance at fixed spp, temporal stability,
light leaking, stale lighting after light changes, runtime, and memory.

\textbf{Expected success.} Unified state reduces variance and temporal shimmer
without requiring a separate cache structure.

\textbf{Failure case.} Radiance and color history need different support sizes
or gating so strongly that they should not be one state.

\subsection{Experiment 3: Motion-Accurate Midpoint From Same State}

\textbf{Purpose.} Test whether the temporal/SR state can also drive frame
interpolation.

\textbf{Scene.} Render $t_0$, $t_{0.5}$, and $t_1$ for:

\begin{itemize}
  \item camera-only motion,
  \item one rigid moving object,
  \item one simple deforming object.
\end{itemize}

\textbf{Method.} Hide $t_{0.5}$. Use Gaussian motion state to synthesize the
midpoint. Compare against the true middle frame.

\textbf{Baselines.}

\begin{itemize}
  \item frame hold,
  \item linear image blend,
  \item screen-space reprojection,
  \item optical-flow interpolation if available.
\end{itemize}

\textbf{Metrics.} Midpoint PSNR/SSIM, silhouette distance, depth-edge alignment,
occlusion error, flow residual, and temporal flicker across $t_0$, generated
$t_{0.5}$, $t_1$.

\textbf{Expected success.} Rigid and camera-motion midpoint silhouettes land
closer to ground truth than screen-space methods.

\textbf{Failure case.} Disoccluded background and non-rigid motion require a
separate residual system, weakening the integration thesis.

\subsection{Experiment 4: Unified State vs Separate State Ablation}

\textbf{Purpose.} Directly answer whether 3DGS integrates features well.

\textbf{Variants.}

\begin{itemize}
  \item image-space-only stack,
  \item separate Gaussian state per feature,
  \item one shared Gaussian state with task-specific subslots,
  \item one shared Gaussian state plus small neural resolver.
\end{itemize}

\textbf{Scorecard.}

\begin{center}
\begin{tabular}{lrrrrr}
\toprule
Variant & TAA & SR & Denoise & Interp & Cost \\
\midrule
Image-space only & ? & ? & ? & ? & ? \\
Separate Gaussians & ? & ? & ? & ? & ? \\
Shared Gaussian state & ? & ? & ? & ? & ? \\
Shared state + neural resolver & ? & ? & ? & ? & ? \\
\bottomrule
\end{tabular}
\end{center}

\textbf{Decision rule.} The shared Gaussian state is a good integration layer
if it wins on at least three of four quality axes while keeping memory/runtime
within roughly 1.5x of the image-space baseline or below the separate-Gaussian
variant.

\textbf{Ablation protocol.} Keep sample inputs, camera jitter, references, and
post-filters identical. Only change the reconstruction state:

\begin{enumerate}
  \item \textbf{Image-space only:} color/depth/velocity history, no Gaussians.
  \item \textbf{Separate Gaussians:} one Gaussian cache for TAA/SR, one for
        radiance, one for interpolation.
  \item \textbf{Shared Gaussians:} one state with task-specific subslots.
  \item \textbf{Shared plus neural:} same shared state, plus equal-size neural
        residual model.
\end{enumerate}

The separate-Gaussian variant is important. If it beats the shared state by a
large margin, then 3DGS may still be useful, but not as the unifying vkGSplat
abstraction.

\subsection{Experiment 5: Downstream Synthetic Data Utility}

\textbf{Purpose.} Test whether reconstruction quality matters for vkGSplat's
synthetic-data goal.

\textbf{Method.} Generate a small robotics/perception dataset with:

\begin{itemize}
  \item noisy low-spp images,
  \item image-space denoised/TAAU images,
  \item unified Gaussian reconstructed images,
  \item high-spp/high-res references.
\end{itemize}

\textbf{Metrics.} Downstream detection/segmentation/pose accuracy, not only
image PSNR.

\textbf{Expected success.} The Gaussian reconstruction improves downstream
accuracy or data efficiency relative to noisy and image-space baselines.

\section{CUDA Validation and Microbenchmarks}

The first implementation should be a CUDA prototype with validation and
microbenchmarks from day one. The tests are wrappers around GPU kernels, not a
CPU research path. A scalar CPU oracle is useful only for tiny deterministic
inputs and expected-value checks.

\begin{itemize}
  \item \textbf{cuda\_extract\_samples:} device kernel unprojects depth,
        primitive ID, color, and motion into compact sample buffers.
  \item \textbf{cuda\_bin\_gaussians:} device hash-grid binning reports cell
        occupancy, overflow rate, atomic pressure, and memory bandwidth.
  \item \textbf{cuda\_assign\_and\_update:} incompatible primitive, depth,
        normal, and motion samples do not update the same Gaussian; benchmark
        branch divergence and update atomics.
  \item \textbf{cuda\_visibility\_decay:} occluded Gaussians decay slowly while
        disoccluded stale Gaussians decay quickly on device.
  \item \textbf{cuda\_project\_taa\_static\_pan:} shared-state TAA reduces
        temporal error on a jittered pan while reporting tile-list reuse.
  \item \textbf{cuda\_project\_sr\_scale2\_4:} the same state resolves 2x and
        4x outputs using only changed projection footprints.
  \item \textbf{cuda\_radiance\_diffuse:} level-1 diffuse radiance cache
        reduces variance against a high-spp Cornell reference without crossing
        primitive boundaries.
  \item \textbf{cuda\_path\_subslot\_moments:} path-depth/lobe moments update
        on device and biased mode is separated from residual/control-variate
        mode.
  \item \textbf{cuda\_midpoint\_rigid:} midpoint interpolation lands a moving
        rigid object's silhouette closer to the rendered middle frame than
        linear image blending.
  \item \textbf{cuda\_framegen\_holes:} disoccluded generated pixels are marked
        as holes and do not receive high-confidence Gaussian color.
  \item \textbf{cuda\_shared\_vs\_separate:} shared-state memory and launch
        count are lower than separate per-feature Gaussian caches while quality
        stays within the scorecard threshold.
\end{itemize}

These benchmarks should use tiny synthetic scenes first, then Wicked/Cornell
captures once the CUDA contracts are stable. Report timing with CUDA events and
memory counters, not only image quality.

\section{Implementation Milestones}

\subsection{M0: CUDA State Without New Rendering}

Start with existing vkGSplat frames and fixtures, but implement the Gaussian
state path as CUDA kernels. Add only device sample extraction, binning,
assignment/update, visibility, and projection. This avoids confusing
reconstruction quality with renderer changes while keeping the research pointed
at the real target: GPU integration.

\textbf{Exit criteria:}

\begin{itemize}
  \item CUDA sample extraction matches the small CPU oracle on tiny inputs,
  \item hash-grid binning has bounded overflow and measurable atomic cost,
  \item visibility decay behaves differently for occlusion and disocclusion,
  \item all kernels expose timing, bandwidth, and debug counters.
\end{itemize}

\subsection{M1: Shared-State TAA}

Implement only $C_j$ and $A_j^{taa}$. Use static camera-pan clips with jitter.

\textbf{Exit criteria:}

\begin{itemize}
  \item temporal error improves over current reprojection,
  \item ghosting does not regress by more than 10\%,
  \item stale history disappears in disoccluded regions within two frames,
  \item TAA projection reuses the same CUDA tile/binning data intended for SR.
\end{itemize}

\subsection{M2: Same State, 2x/4x SR}

Add $A_j^{sr}$ but keep the same shared core $C_j$. The only scale-specific
change should be the projection footprint.

\textbf{Exit criteria:}

\begin{itemize}
  \item 2x SR beats bilinear and screen-space TAAU on thin geometry,
  \item 4x SR does not require a second Gaussian cache,
  \item blur is bounded by an edge-width metric against high-res reference,
  \item CUDA projection cost scales with output pixels and tile occupancy, not
        by rebuilding the Gaussian state.
\end{itemize}

\subsection{M3: Radiance Subslot}

Add $A_j^{rad}$ for level-1 diffuse/rough indirect radiance. Keep direct light
sampled normally.

\textbf{Exit criteria:}

\begin{itemize}
  \item variance drops at fixed spp,
  \item light leaking across primitive/depth boundaries stays below baseline
        denoiser artifacts,
  \item TAA/SR quality does not degrade when radiance subslots are enabled,
  \item radiance subslot updates do not dominate frame time through atomics.
\end{itemize}

\subsection{M4: Motion Subslot}

Add $A_j^{mot}$ and synthesize true midpoint frames for rigid object and
camera-only motion.

\textbf{Exit criteria:}

\begin{itemize}
  \item silhouette IoU beats image blending and screen reprojection,
  \item flow residual predicts failures,
  \item disoccluded background is masked rather than hallucinated as confident,
  \item generated frames reuse Gaussian motion/tile data rather than launching
        an unrelated image-space interpolation stack.
\end{itemize}

\subsection{M5: Shared State Plus Neural Residual}

Only after M1--M4 show signal, train a small equal-size resolver with and
without Gaussian feature planes.

\textbf{Exit criteria:}

\begin{itemize}
  \item Gaussian features improve error at fixed model size,
  \item model does not simply learn to ignore $\Phi_t$,
  \item inference remains a small residual pass, not the whole system.
\end{itemize}

\section{CUDA Implementation Sketch}

Start with CUDA-facing contracts. Host structs define buffer layout and launch
configuration; the persistent Gaussian state lives in device structure-of-arrays
buffers:

\begin{verbatim}
struct GaussianSampleDeviceView {
    float3* position;
    float3* normal;
    float3* color;
    float3* radiance;
    float3* velocity;
    float2* motion_px;
    float* depth;
    float* variance;
    float* roughness;
    uint32_t* primitive_id;
    uint32_t* material_id;
    uint32_t* path_depth;
    uint32_t count;
};

struct PathLobeMoments {
    float3 mean;
    float second_moment;
    float sample_weight;
    float robust_residual;
};

struct GaussianStateDeviceView {
    float3* mean;
    float4* covariance_0_3;
    float2* covariance_4_5;
    float4* rotation;
    float3* color;
    float3* radiance_levels[3];
    PathLobeMoments* path_lobes;
    float3* velocity;
    float* variance;
    float* confidence;
    float* visibility_age;
    float* flow_residual_px;
    uint32_t* primitive_id;
    uint32_t* material_id;
    uint32_t* motion_group_id;
    uint32_t count;
};
\end{verbatim}

CUDA launch API:

\begin{verbatim}
class CudaGaussianReconstructionCache {
public:
    void extract_samples_cuda(const SeedFrameDeviceView& seed);
    void update_visibility_cuda(const DepthIdDeviceView& current);
    void bin_gaussians_cuda();
    void assign_and_update_cuda();
    void build_tile_lists_cuda(const CameraDeviceConstants& camera);
    void project_taa_cuda(ReconstructionFrameDeviceView out) const;
    void project_super_resolution_cuda(ReconstructionFrameDeviceView out) const;
    void query_indirect_cuda(ShadingQueryDeviceView queries) const;
    void update_path_subslots_cuda(PathEventDeviceView events);
    void synthesize_generated_cuda(float alpha, FrameDeviceView out) const;
    void project_features_cuda(FeatureFrameDeviceView out) const;
};
\end{verbatim}

The first public C++ API can wrap these launches, but the design should be
judged by CUDA memory layout, launch count, and tile-list reuse.

\section{Numerical Invariants and Debug Counters}

The first CUDA implementation should treat every equation above as an invariant
that can fail loudly. This makes the research question falsifiable inside the
GPU path, not after translating a CPU prototype.

\subsection{Hard Invariants}

\begin{enumerate}
  \item \textbf{Confidence is bounded:}
  \[
    0 \leq \nu_j \leq 1.
  \]
  Any update that produces a value outside this range is a bug, not a tuning
  problem.

  \item \textbf{Covariance stays positive definite:}
  \[
    \lambda_{min} I \preceq \Sigma_j \preceq \lambda_{max} I.
  \]
  Clamp eigenvalues in the CUDA update kernel after every update. A clamped
  covariance should increment a device debug counter; frequent clamps mean
  assignment or split rules are unstable.

  \item \textbf{Primitive identity is a hard early gate:}
  \[
    id_i^{prim}\neq id_j^{prim}
    \quad\Rightarrow\quad
    W_{ij}=0
  \]
  when both IDs are valid. Later experiments can soften this for alpha-tested
  or deforming surfaces, but the first demo should prove that 3DGS history does
  not cross triangle boundaries.

  \item \textbf{Disoccluded history is never high confidence:}
  \[
    D_p=1 \Rightarrow \sum_{j\rightarrow p} a_{jp}\nu_j \leq \epsilon_D.
  \]
  This is the main anti-ghosting invariant. If a newly revealed pixel receives
  strong old Gaussian history, the shared-state thesis fails for TAA.

  \item \textbf{Radiance estimates stay non-negative and bounded:}
  \[
    0 \leq \hat{L}_j^k \leq L_{max,k}.
  \]
  The bound is not a physically pure estimator; it is a prototype guard against
  rare path samples poisoning TAA/SR history.

  \item \textbf{Stable assignment:} if two Gaussians have equal compatibility,
  the lower stable index wins. This keeps CUDA validation reproducible enough
  to compare runs while accepting normal floating-point ordering differences.
\end{enumerate}

\subsection{Per-Frame Counters}

Every run should report:

\begin{itemize}
  \item samples extracted, assigned, rejected, and spawned,
  \item rejections by gate: primitive, depth, normal, motion, residual,
        confidence,
  \item covariance clamps, splits, merges, and prunes,
  \item visible, occluded, disoccluded, and out-of-frustum Gaussian counts,
  \item stale-history energy on disoccluded pixels,
  \item memory bytes for shared state versus separate-feature states,
  \item per-feature quality deltas: TAA, SR, radiance, interpolation.
\end{itemize}

The key diagnostic is stale-history energy:

\[
  E_{stale} =
  \frac{1}{|D_t|}
  \sum_{p:D_t(p)=1}
  \left\|
    \sum_{j\rightarrow p} a_{jp}\nu_j c_j
  \right\|_1.
\]

If $E_{stale}$ stays high after tuning gates, world-space Gaussians are not a
safe TAA history representation for the current scene class.

\section{CUDA Performance Targets}

Quality wins are not enough. A shared Gaussian state is valuable only if the
CUDA execution graph is simpler or cheaper than separate feature stacks.

\subsection{Primary GPU Metrics}

Each experiment should report:

\begin{itemize}
  \item per-kernel time from CUDA events,
  \item end-to-end reconstruction frame time,
  \item global memory read/write bytes,
  \item atomics per sample and atomics per Gaussian,
  \item hash-cell overflow rate,
  \item tile-list build cost and tile-list reuse count,
  \item average and worst-case Gaussians per tile,
  \item occupancy, register pressure, and shared-memory use for update and
        projection kernels,
  \item memory footprint for image-space baselines, separate Gaussian states,
        and the shared Gaussian state.
\end{itemize}

\subsection{GPU Integration Score}

The CUDA-specific score should separate visual quality from implementation
cost:

\[
  S_{cuda} =
  S_{quality}
  - \lambda_t \frac{T_{shared}}{T_{image}}
  - \lambda_m \frac{M_{shared}}{M_{image}}
  + \lambda_r R_{reuse}.
\]

$R_{reuse}$ measures how often the same Gaussian bins, tile lists, confidence
values, and visibility decisions feed multiple outputs:

\[
  R_{reuse} =
  \frac{
    N_{shared\_consumers}
  }{
    N_{extra\_builds}+1
  }.
\]

The shared-state thesis is weak if $S_{quality}>0$ but $S_{cuda}<0$. That case
means the idea may be mathematically interesting but not a good vkGSplat CUDA
architecture.

\subsection{CUDA Kill Conditions}

Stop or redesign the architecture if:

\begin{itemize}
  \item assignment requires unbounded neighbor search,
  \item atomics dominate update time before scenes become interesting,
  \item TAA, SR, radiance, and frame generation each rebuild their own bins,
  \item tile occupancy variance creates persistent long-tail frame times,
  \item shared-state memory exceeds separate Gaussian states,
  \item quality depends on a CPU-only ordering or exact determinism assumption.
\end{itemize}

\section{Parameter Sweeps and Thesis Kill Tests}

The first implementation should not hand-tune one pleasant clip. It should run
small sweeps over the parameters that decide whether the unified state is
actually robust.

\begin{center}
\begin{tabular}{llll}
\toprule
Parameter & Sweep & Primary metric & Kill condition \\
\midrule
$\theta_{assign}$ & 0.05--0.60 & $E_{ghost},E_{temp}$ & no stable interval \\
$\sigma_d$ & 0.002--0.08 & disocclusion error & depth leak unavoidable \\
$\sigma_n$ & 0.05--0.50 & edge shimmer & normal gate unused \\
$\lambda_{min}$ & $10^{-6}$--$10^{-3}$ & blur, clamps & covariance collapses \\
$\rho_{SR}$ & 0.35--1.00 & edge width & 2x requires new state \\
$\gamma_{dis}$ & 0.05--0.50 & stale energy & stale energy persists \\
$K_{cell}$ & 8--64 & memory, quality & quality needs unbounded K \\
$\eta_{\dot{\mu}}$ & 0.05--0.80 & $E_{mid},E_{sil}$ & residual cannot predict fail \\
\bottomrule
\end{tabular}
\end{center}

For each sweep point, compute a compact integration score:

\[
  S =
  w_T \Delta E_{temp} +
  w_R \Delta E_{ref}^{SR} +
  w_L \Delta V_L +
  w_M \Delta E_{mid} -
  w_G \max(0,\Delta E_{ghost}) -
  w_C \Delta C_{cost}.
\]

Here positive $\Delta$ means improvement over the image-space baseline, except
for ghosting and cost where increases are penalties. The unified state passes
the research bar only if there is a non-tiny parameter interval where $S>0$ and
$E_{ghost}$ remains within the 10\% guardrail. A single winning parameter point
is not enough.

\subsection{Ablation Protocol}

The decisive ablation should hold sample inputs, jitter, and camera path fixed:

\begin{enumerate}
  \item \textbf{Image-space only:} current reprojection plus SVGF-style denoise.
  \item \textbf{Separate Gaussian states:} one Gaussian cache for TAA, one for
        SR, one for radiance, one for motion.
  \item \textbf{Shared core:} one $C_j$ with task subslots.
  \item \textbf{Shared core plus neural residual:} same as above, with
        projected Gaussian feature planes.
\end{enumerate}

The integration hypothesis is supported only if variant 3 beats variant 2 on
cost or consistency while staying close in quality. Variant 4 is useful only if
the learned residual uses Gaussian features; test this by zeroing $\Phi_t$ at
validation time and measuring the error jump.

\section{Minimal Demonstration in the Current vkGSplat Workspace}

The lowest-risk implementation path fits the current repository shape:

\begin{itemize}
  \item add \texttt{include/vkgsplat/cuda/gaussian\_reconstruction.h},
  \item add \texttt{src/cuda/gaussian\_reconstruction.cu},
  \item add \texttt{tests/test\_cuda\_gaussian\_reconstruction.cu},
  \item register the CUDA benchmark/test only under
        \texttt{VKGSPLAT\_ENABLE\_CUDA}.
\end{itemize}

The first CUDA code should reuse existing frame contracts instead of inventing
new scene loading:

\begin{itemize}
  \item \texttt{RayTracingSeedFrame} supplies color, depth, NDC depth,
        primitive ID, and camera matrices,
  \item \texttt{compute\_camera\_motion\_map} supplies the baseline motion map,
  \item \texttt{reproject\_history} supplies the image-space TAA baseline,
  \item \texttt{denoise\_svgf\_baseline} supplies the Monte Carlo denoise
        baseline.
\end{itemize}

The first demo sequence should be tiny:

\begin{enumerate}
  \item render four low-resolution jittered seed frames from the existing
        triangle/Cornell-style ray-tracing seed path,
  \item upload seed frames and auxiliary buffers to device memory,
  \item build the Gaussian cache from frames 0--2 using CUDA update kernels,
  \item project frame 3 at native resolution for TAA comparison,
  \item project frame 3 at 2x resolution for SR comparison,
  \item synthesize the midpoint between frames 1 and 2 on device, comparing
        against a held-out rendered middle frame,
  \item optionally query level-1 radiance from the same Gaussian state for a
        low-spp diffuse indirect test,
  \item report CUDA launch time, memory traffic, atomics, tile occupancy, and
        quality metrics for every run.
\end{enumerate}

This demo proves integration only if the same Gaussian indices and shared core
fields are used across all projections. A convenient implementation check is:

\[
  \frac{|J_{taa}\cap J_{sr}\cap J_{mot}|}
       {|J_{taa}\cup J_{sr}\cup J_{mot}|}
  \geq 0.6
\]

where $J_{task}$ is the set of Gaussian IDs that materially contribute to a
task output. If overlap is low, the method is drifting toward separate hidden
systems.

\section{DLSS-5-Style Control Plane}

If the analytic state works, the later neural resolver should not receive only
RGB history. It should receive a control plane that makes generated detail
accountable to the explicit scene state:

\[
  X_t =
  (C_t, Z_t, U_t, N_t, A_t,
   \Phi_t^{gauss}, \nu_t^{gauss}, V_t^{gauss}, D_t).
\]

Here $C_t$ is low-resolution color, $Z_t$ depth, $U_t$ motion, $N_t$ normal,
$A_t$ material/albedo metadata, $\Phi_t^{gauss}$ projected Gaussian features,
$\nu_t^{gauss}$ confidence, $V_t^{gauss}$ variance, and $D_t$ disocclusion.
The resolver output should be decomposed:

\[
  O_t = O_t^{analytic} + R_\theta(X_t),
\]

where $O_t^{analytic}$ is the deterministic Gaussian reconstruction and
$R_\theta$ is a bounded residual. This keeps the DLSS-5-style part honest: the
network can add high-frequency lighting/material corrections, but the analytic
state remains responsible for geometry, motion, visibility, and uncertainty.

\section{Risks}

\begin{itemize}
  \item One state may be too overloaded; TAA, radiance, and interpolation may
        need incompatible covariance scales.
  \item Ghosting may worsen if stale Gaussians persist after disocclusion.
  \item Radiance caching can leak light across visibility boundaries.
  \item Super-resolution can become stable but blurry.
  \item Frame interpolation needs residual handling for particles, UI,
        transparency, reflections, and animated shading.
  \item A DLSS-class commercial product still needs large training,
        validation, and latency engineering beyond the explicit Gaussian
        state, but the core vkGSplat framework should not depend on that.
\end{itemize}

\section{Current Recommendation}

Proceed with the unified 3DGS integration thesis. The first target should not
be a DLSS clone. The first target should be a falsifiable demo:

\begin{quote}
One Gaussian state, built from low-resolution jittered seed frames, drives both
TAA and 2x super-resolution on a thin-geometry camera-pan clip, then reuses the
same state for a rigid-motion midpoint interpolation test.
\end{quote}

If that works, add radiance-cache subslots. Only after the no-training path has
a measurable signal should vkGSplat add a small learned resolver. If it does
not, 3DGS may still be useful for individual features, but not as the
integrating idea behind vkGSplat.

\appendix

\section{Appendix: Starting from Zero Graphics Hardware}

If a vkGSplat-oriented accelerator starts from zero graphics-specific hardware,
it should not attempt to recreate a full conventional GPU pipeline first. The
right target is a neural-rendering/GS accelerator with a minimal graphics-like
prefill front end. Graphics hardware should be treated as an evidence generator,
not as the center of the product.

The minimum useful prefill hardware is:

\begin{itemize}
  \item tiled triangle or mesh rasterization for visibility, depth, primitive
        identity, and barycentrics,
  \item depth test and preferably hierarchical/tile Z,
  \item barycentric attribute interpolation for position, normal, material, and
        motion features,
  \item basic texture sampling only if material evidence cannot be supplied
        another way,
  \item simple compositing/blending for seed buffers,
  \item optional BVH traversal only if ray evidence is central to the product.
\end{itemize}

The blocks to avoid in a first architecture are the broad legacy pieces that
exist mostly to support full graphics API completeness:

\begin{itemize}
  \item a complex fixed-function shader pipeline,
  \item tessellation, geometry-shader, and legacy compatibility paths,
  \item a complete texture/filtering ecosystem beyond the needed evidence
        channels,
  \item a large ROP/blend feature matrix,
  \item heavyweight generalized graphics scheduling,
  \item large RT traversal blocks unless radiance/path evidence is a primary
        workload.
\end{itemize}

The evidence target is not a finished shaded frame. The target is a tensor
bundle:

\[
  A_t=(L,D,ID,X,N,V,M,\nu,\sigma^2),
\]

where $L$ is color/radiance evidence, $D$ is depth, $ID$ is primitive/material
identity, $X$ is world position, $N$ is normal, $V$ is motion, $M$ is material
or mask data, $\nu$ is confidence, and $\sigma^2$ is variance. Once this bundle
is trustworthy, the differentiating work moves into Gaussian update,
radiance-cache paging, tensorized decode, and optional neural resolving.

For a first vkGSplat-class architecture, prioritize hardware roughly as:

\begin{center}
\small
\begin{tabular}{cp{0.34\linewidth}p{0.44\linewidth}}
\toprule
Priority & Hardware block & Reason \\
\midrule
1 & tensor/matrix engines & Gaussian feature projection, attention-like resolve, and neural decode \\
2 & SRAM/cache bandwidth and paging & radiance/KV-like cache residency and streaming reuse \\
3 & sparse gather/scatter, reductions, atomics & online Gaussian update, compaction, assignment, and binning \\
4 & tile scheduler and block/page tables & local decode, work clustering, and cache management \\
5 & minimal raster/depth/interpolation & trustworthy prefill evidence for visibility, IDs, depth, and motion \\
6 & texture sampling & only where material evidence cannot be represented cheaply as attributes \\
7 & BVH/ray traversal & only if path/radiance evidence is central to the product target \\
\bottomrule
\end{tabular}
\end{center}

The architectural pipeline should be:

\[
\begin{gathered}
\text{minimal raster/ray prefill}
\rightarrow \text{seed evidence tensors}
\rightarrow \text{Gaussian update}\\
\rightarrow \text{paged radiance cache}
\rightarrow \text{tensor decode for TAA/SR/framegen}.
\end{gathered}
\]

This is different from dropping graphics entirely. The accelerator still needs
visibility, identity, depth, and material evidence. But it does not need to
spend early area on a general-purpose graphics pipeline if the product thesis
is persistent Gaussian state and DLSS-class decode. The design principle is:

\begin{quote}
Do not reinvent graphics hardware wholesale. Reinvent the boundary: a thin
prefill evidence generator feeding a tensorized Gaussian/radiance-cache decode
engine.
\end{quote}

A useful pass/fail test for hardware planning is to implement two simulators:

\begin{enumerate}
  \item \textbf{Minimal evidence prefill:} raster/depth/ID/interpolation only,
        plus tensor decode.
  \item \textbf{Richer traditional prefill:} add texture/filtering/ray features
        progressively.
\end{enumerate}

If richer prefill improves reconstruction quality much more than it costs, add
those graphics blocks. If the minimal prefill already gives stable Gaussian
state and good decode quality, invest the saved area in tensor/cache/paging
instead. The burden of proof should be on adding legacy graphics area, not on
removing the core tensor/cache path.

\section{Appendix: Architectural Path Dependence}

The deeper vkGSplat thesis is not only that 3DGS can run on tensor cores. It is
that computer architecture is shaped by dominant paradigms, and then the
resulting programmable substrate becomes available to later paradigms. A
successful workload creates hardware pressure; hardware vendors respond with
specialized but still programmable mechanisms; a different representation later
discovers that those mechanisms are close to the right machine for its own
needs.

A compact form of the pattern is:

\[
\begin{gathered}
\text{dominant paradigm}
\rightarrow \text{hardware specialization}\\
\rightarrow \text{programmable affordance}
\rightarrow \text{new paradigm reuse}.
\end{gathered}
\]

Historically, raster graphics did this first. The need to draw shaded triangles
at high throughput justified programmable GPU shaders, high memory bandwidth,
texture units, and massive data parallelism. Once exposed as a programmable
substrate, that graphics-shaped machine became useful for scientific computing,
deep learning, differentiable rendering, NeRFs, and other non-raster workloads.
The hardware was motivated by graphics, but it did not remain confined to
classical graphics.

Transformer/LLM serving is now applying the same pressure in a different
region of the design space. It justifies tensor engines, FP8/FP4 arithmetic,
block-scaled quantization, KV-cache paging, high-bandwidth memory, large SRAM,
continuous batching, and decode-oriented scheduling \cite{nvidia2022hopper,
nvidia2026transformerenginefp8,nvidia2026nvfp4,kwon2023pagedattention,
prabhu2025vattention}. These features are not graphics-specific, but they are
surprisingly well matched to Gaussian-state reconstruction: tile-local feature
projection, paged radiance caches, motion-aligned history, low-precision
feature packets, sparse gather/reduce, and optional transformer decoding.

\begin{center}
\small
\begin{tabular}{p{0.24\linewidth}p{0.32\linewidth}p{0.33\linewidth}}
\toprule
Dominant paradigm & Architecture it encourages & Later reuse \\
\midrule
Raster graphics & Programmable shaders, high bandwidth, texture access, SIMT execution & GPGPU, deep learning, differentiable rendering, neural fields \\
Deep learning and LLMs & Tensor cores, FP8/FP4, HBM, SRAM, cache paging, batched decode & 3DGS feature projection, radiance cache, TAA/SR/framegen decode \\
Stateful 3DGS reconstruction & Gaussian/tile pages, sparse updates, visibility-aware caches, narrow evidence prefill & Future digital twins, robotics SDG, neural-rendering accelerators \\
\bottomrule
\end{tabular}
\end{center}

For vkGSplat, the inversion is important. Graphics hardware helped create the
GPU. AI then reshaped the GPU around tensor and cache workloads. Now a new
graphics representation, 3D Gaussian Splatting, may benefit more from the
AI-shaped architecture than from the older fixed-function graphics pipeline.
3DGS is therefore not competing with LLM hardware. It is discovering that LLM
hardware has already built much of the right substrate for stateful graphics
reconstruction.

This view also explains why vkGSplat should not begin by cloning a full legacy
GPU. The prefill side still needs graphics-like evidence: visibility, depth,
primitive identity, barycentrics, material hints, normals, and motion. But the
decode side is where the architectural surplus lies. Gaussian state update,
radiance-cache lookup, projected feature tokens, temporal reconstruction,
super-resolution, and frame generation all map to the tensor/cache/paging
machinery that LLM serving made economically inevitable.

The appendix-level claim is therefore:

\begin{quote}
vkGSplat is an example of architectural path dependence. Transformer/LLM
serving justifies hardware investment in tensor engines, low-precision
arithmetic, cache paging, and high-throughput decode scheduling. Those
investments then become programmable affordances for a different paradigm:
Gaussian-state graphics reconstruction.
\end{quote}

The falsifiable version is the same as the hardware ablation in the main text.
If a minimal evidence prefill plus strong tensor/cache decode beats a
richer graphics prefill under equal area, bandwidth, or energy budgets, then
vkGSplat is benefiting from LLM-shaped hardware. If richer prefill dominates,
then the old graphics paradigm still owns the bottleneck. The interesting
outcome is not ideological; it is the measured boundary where one dominant
paradigm's hardware surplus becomes another paradigm's accelerator.

\begin{thebibliography}{99}

\bibitem{kerbl2023gs}
Bernhard Kerbl, Georgios Kopanas, Thomas Leimkuehler, and George Drettakis.
\newblock 3D Gaussian Splatting for Real-Time Radiance Field Rendering.
\newblock ACM Transactions on Graphics, 2023.
\newblock \url{https://repo-sam.inria.fr/fungraph/3d-gaussian-splatting/}

\bibitem{laine2020nvdiffrast}
Samuli Laine, Janne Hellsten, Tero Karras, Yeongho Seol, Jaakko Lehtinen,
and Timo Aila.
\newblock Modular Primitives for High-Performance Differentiable Rendering.
\newblock ACM Transactions on Graphics, 2020.
\newblock \url{https://nvlabs.github.io/nvdiffrast/}

\bibitem{ye2024gsplat}
Vickie Ye, Ruilong Li, Justin Kerr, Angjoo Kanazawa, and Matthew Tancik.
\newblock gsplat: An Open-Source Library for Gaussian Splatting.
\newblock arXiv:2409.06765, 2024.
\newblock \url{https://arxiv.org/abs/2409.06765}

\bibitem{yu2024mipsplat}
Zehao Yu, Anpei Chen, Binbin Huang, Torsten Sattler, and Andreas Geiger.
\newblock Mip-Splatting: Alias-Free 3D Gaussian Splatting.
\newblock CVPR, 2024.
\newblock \url{https://openaccess.thecvf.com/content/CVPR2024/html/Yu_Mip-Splatting_Alias-free_3D_Gaussian_Splatting_CVPR_2024_paper.html}

\bibitem{yan2024multiscale}
Zhiwen Yan, Weng Fei Low, Yu Chen, and Gim Hee Lee.
\newblock Multi-Scale 3D Gaussian Splatting for Anti-Aliased Rendering.
\newblock CVPR, 2024.
\newblock \url{https://cvpr.thecvf.com/virtual/2024/poster/30347}

\bibitem{liang2024analytic}
Zhihao Liang, Qi Zhang, Wenbo Hu, Ying Feng, Lei Zhu, and Kui Jia.
\newblock Analytic-Splatting: Anti-Aliased 3D Gaussian Splatting via Analytic
Integration.
\newblock arXiv:2403.11056, 2024.
\newblock \url{https://arxiv.org/abs/2403.11056}

\bibitem{steiner2025aaa}
Michael Steiner, Thomas Koehler, Lukas Radl, Felix Windisch,
Dieter Schmalstieg, and Markus Steinberger.
\newblock AAA-Gaussians: Anti-Aliased and Artifact-Free 3D Gaussian Rendering.
\newblock arXiv preprint, 2025.
\newblock \url{https://tugraz.elsevierpure.com/en/publications/aaa-gaussians-anti-aliased-and-artifact-free-3d-gaussian-renderin/}

\bibitem{yang2025lodgs}
Zhenya Yang, Bingchen Gong, and Kai Chen.
\newblock LOD-GS: Level-of-Detail-Sensitive 3D Gaussian Splatting for Detail
Conserved Anti-Aliasing.
\newblock arXiv:2507.00554, 2025.
\newblock \url{https://papers.cool/arxiv/2507.00554}

\bibitem{suh2026aasplat}
Taewoo Suh, Sungpyo Kim, Jongmin Park, and Munchurl Kim.
\newblock AA-Splat: Anti-Aliased Feed-forward Gaussian Splatting.
\newblock arXiv:2603.29394, 2026.
\newblock \url{https://papers.cool/arxiv/2603.29394}

\bibitem{bauer2026gscache}
David Bauer, Qi Wu, Hamid Gadirov, and Kwan Liu Ma.
\newblock GSCache: Real-Time Radiance Caching for Volume Path Tracing Using
3D Gaussian Splatting.
\newblock IEEE Transactions on Visualization and Computer Graphics, 2026.
\newblock \url{https://research.rug.nl/en/publications/gscache-real-time-radiance-caching-for-volume-path-tracing-using-/}

\bibitem{muller2021nrc}
Thomas Mueller, Fabrice Rousselle, Jan Novak, and Alexander Keller.
\newblock Real-Time Neural Radiance Caching for Path Tracing.
\newblock ACM Transactions on Graphics, 2021.
\newblock \url{https://research.nvidia.com/labs/rtr/publication/muller2021nrc/}

\bibitem{muller2020ncv}
Thomas Mueller, Fabrice Rousselle, Alexander Keller, and Jan Novak.
\newblock Neural Control Variates.
\newblock ACM Transactions on Graphics, 2020.
\newblock \url{https://research.nvidia.com/labs/rtr/publication/muller2020ncv/}

\bibitem{schied2017svgf}
Christoph Schied, Anton Kaplanyan, Chris Wyman, Anjul Patney,
Chakravarty R. Alla Chaitanya, John Burgess, Shiqiu Liu,
Carsten Dachsbacher, Aaron Lefohn, and Marco Salvi.
\newblock Spatiotemporal Variance-Guided Filtering: Real-Time Reconstruction
for Path-Traced Global Illumination.
\newblock High Performance Graphics, 2017.
\newblock \url{https://research.nvidia.com/labs/rtr/publication/schied2017spatiotemporal/}

\bibitem{bako2017kpcn}
Steve Bako, Thijs Vogels, Brian McWilliams, Mark Meyer, Jan Novak,
Alex Harvill, Pradeep Sen, Tony DeRose, and Fabrice Rousselle.
\newblock Kernel-Predicting Convolutional Networks for Denoising Monte Carlo
Renderings.
\newblock ACM Transactions on Graphics, 2017.
\newblock \url{https://studios.disneyresearch.com/2017/07/20/kernel-predicting-convolutional-networks-for-denoising-monte-carlo-renderings/}

\bibitem{bitterli2020restir}
Benedikt Bitterli, Chris Wyman, Matt Pharr, Peter Shirley, Aaron Lefohn,
and Wojciech Jarosz.
\newblock Spatiotemporal Reservoir Resampling for Real-Time Ray Tracing with
Dynamic Direct Lighting.
\newblock ACM Transactions on Graphics, 2020.
\newblock \url{https://research.nvidia.com/labs/rtr/publication/bitterli2020spatiotemporal/}

\bibitem{kheradmand2025stochasticsplats}
Shakiba Kheradmand, Delio Vicini, George Kopanas, Dmitry Lagun,
Kwang Moo Yi, Mark Matthews, and Andrea Tagliasacchi.
\newblock StochasticSplats: Stochastic Rasterization for Sorting-Free 3D
Gaussian Splatting.
\newblock arXiv:2503.24366, 2025.
\newblock \url{https://arxiv.org/abs/2503.24366}

\bibitem{zhou2024unified}
Yang Zhou, Songyin Wu, and Ling-Qi Yan.
\newblock Unified Gaussian Primitives for Scene Representation and Rendering.
\newblock arXiv:2406.09733, 2024.
\newblock \url{https://mangosister.github.io/gsr_site/}

\bibitem{wu2024dynamic4dgs}
Guanjun Wu, Taoran Yi, Jiemin Fang, Lingxi Xie, Xiaopeng Zhang, Wei Wei,
Wenyu Liu, Qi Tian, and Xinggang Wang.
\newblock 4D Gaussian Splatting for Real-Time Dynamic Scene Rendering.
\newblock CVPR, 2024.
\newblock \url{https://guanjunwu.github.io/4dgs/}

\bibitem{li2023spacetime}
Zhan Li, Zhang Chen, Zhong Li, and Yi Xu.
\newblock Spacetime Gaussian Feature Splatting for Real-Time Dynamic View
Synthesis.
\newblock arXiv:2312.16812, 2023.
\newblock \url{https://arxiv.org/abs/2312.16812}

\bibitem{duan2024rotor4dgs}
Yuanxing Duan, Fangyin Wei, Qiyu Dai, Yuhang He, Wenzheng Chen, and Baoquan Chen.
\newblock 4D-Rotor Gaussian Splatting: Towards Efficient Novel View Synthesis
for Dynamic Scenes.
\newblock 2024.
\newblock \url{https://research.nvidia.com/publication/2024-02_4d-rotor-gaussian-splatting-towards-efficient-novel-view-synthesis-dynamic}

\bibitem{lin2023gaussianflow}
Youtian Lin, Zuozhuo Dai, Siyu Zhu, and Yao Yao.
\newblock Gaussian-Flow: 4D Reconstruction with Dynamic 3D Gaussian Particle.
\newblock arXiv:2312.03431, 2023.
\newblock \url{https://huggingface.co/papers/2312.03431}

\bibitem{gao2024gaussianflow}
Quankai Gao, Qiangeng Xu, Zhe Cao, Ben Mildenhall, Wenchao Ma, Le Chen,
Dahang Tang, and Ulrich Neumann.
\newblock GaussianFlow: Splatting Gaussian Dynamics for 4D Content Creation.
\newblock arXiv:2403.12365, 2024.
\newblock \url{https://huggingface.co/papers/2403.12365}

\bibitem{zhu2024motiongs}
Ruijie Zhu, Yanzhe Liang, Hanzhi Chang, Jiacheng Deng, Jiahao Lu,
Wenfei Yang, Tianzhu Zhang, and Yongdong Zhang.
\newblock MotionGS: Exploring Explicit Motion Guidance for Deformable 3D
Gaussian Splatting.
\newblock arXiv:2410.07707, 2024.
\newblock \url{https://huggingface.co/papers/2410.07707}

\bibitem{liu2026kfgs}
Qingyuan Tang, Yufei Yin, Yanming Zhu, Zhou Yu, Zhenzhong Kuang,
Jiajun Ding, and Jifa He.
\newblock KF-GS: Kalman Filter-Guided Gaussian Splatting for Real-Time
High-Quality Dynamic Scene Reconstruction.
\newblock Journal of Visual Communication and Image Representation, in press,
2026.
\newblock \url{https://www.sciencedirect.com/science/article/pii/S1047320326001100}

\bibitem{wang2026retimegs}
Xuezhen Wang, Li Ma, Yulin Shen, Zeyu Wang, and Pedro V. Sander.
\newblock RetimeGS: Continuous-Time Reconstruction of 4D Gaussian Splatting.
\newblock CVPR, 2026.
\newblock \url{https://william-wang2.github.io/RetimeGS/}

\bibitem{nvidia2026dlss5}
NVIDIA.
\newblock NVIDIA DLSS 5 Delivers AI-Powered Breakthrough in Visual Fidelity
for Games.
\newblock March 16, 2026.
\newblock \url{https://nvidianews.nvidia.com/news/nvidia-dlss-5-delivers-ai-powered-breakthrough-in-visual-fidelity-for-games}

\bibitem{nvidia2026dlss45}
NVIDIA.
\newblock DLSS 4.5 Dynamic Multi Frame Generation and Multi Frame Generation
6X Available Now.
\newblock March 31, 2026.
\newblock \url{https://www.nvidia.com/en-us/geforce/news/dlss-4-5-dynamic-multi-frame-generation-6x-mode-released/}

\bibitem{nvidia2026dlssdeveloper}
NVIDIA Developer.
\newblock NVIDIA DLSS.
\newblock Accessed May 5, 2026.
\newblock \url{https://developer.nvidia.com/rtx/dlss}

\bibitem{nvidia2025dlss4}
NVIDIA.
\newblock NVIDIA DLSS 4 Introduces Multi Frame Generation and Enhancements for
All DLSS Technologies.
\newblock January 6, 2025.
\newblock \url{https://www.nvidia.com/en-sg/geforce/news/dlss4-multi-frame-generation-ai-innovations/}

\bibitem{nvidia2025dlss4research}
NVIDIA Research.
\newblock DLSS 4: Transforming Real-Time Graphics with AI.
\newblock Technical report, 2025.
\newblock \url{https://research.nvidia.com/labs/adlr/DLSS4/}

\bibitem{nvidia2022hopper}
NVIDIA.
\newblock NVIDIA Hopper Architecture In-Depth.
\newblock Technical blog, 2022.
\newblock \url{https://developer.nvidia.com/blog/nvidia-hopper-architecture-in-depth/}

\bibitem{nvidia2026transformerenginefp8}
NVIDIA.
\newblock FP8 Current Scaling.
\newblock Transformer Engine documentation, accessed May 5, 2026.
\newblock \url{https://docs.nvidia.com/deeplearning/transformer-engine/user-guide/features/low_precision_training/fp8_current_scaling/fp8_current_scaling.html}

\bibitem{nvidia2026cudamath}
NVIDIA.
\newblock CUDA Math API Reference Manual.
\newblock CUDA Toolkit documentation, accessed May 5, 2026.
\newblock \url{https://docs.nvidia.com/cuda/archive/13.0.1/cuda-math-api/index.html}

\bibitem{nvidia2026nvfp4}
NVIDIA.
\newblock NVFP4.
\newblock Transformer Engine documentation, accessed May 5, 2026.
\newblock \url{https://nvidia.github.io/TransformerEngine/features/low_precision_training/nvfp4/nvfp4.html}

\bibitem{ocp2026mx}
Open Compute Project.
\newblock OCP Microscaling Formats (MX) Specification, Version 1.0.
\newblock 2026.
\newblock \url{https://www.opencompute.org/documents/ocp-microscaling-formats-mx-v1-0-spec-final-pdf}

\bibitem{kwon2023pagedattention}
Woosuk Kwon, Zhuohan Li, Siyuan Zhuang, Ying Sheng, Lianmin Zheng,
Cody Hao Yu, Joseph E. Gonzalez, Hao Zhang, and Ion Stoica.
\newblock Efficient Memory Management for Large Language Model Serving with
PagedAttention.
\newblock SOSP, 2023.
\newblock \url{https://huggingface.co/papers/2309.06180}

\bibitem{prabhu2025vattention}
Ramya Prabhu, Ajay Nayak, Jayashree Mohan, Ramachandran Ramjee, and
Ashish Panwar.
\newblock vAttention: Dynamic Memory Management for Serving LLMs without
PagedAttention.
\newblock ASPLOS, 2025.
\newblock \url{https://www.microsoft.com/en-us/research/publication/vattention-dynamic-memory-management-for-serving-llms-without-pagedattention/}

\bibitem{intel2026xess}
Intel.
\newblock Intel XeSS Super Resolution Developer Guide 2.0.
\newblock Accessed May 5, 2026.
\newblock \url{https://www.intel.com/content/www/us/en/developer/articles/technical/xess-sr-developer-guide.html}

\bibitem{amd2026fsr}
AMD.
\newblock AMD FSR Technologies.
\newblock Accessed May 5, 2026.
\newblock \url{https://www.amd.com/en/products/graphics/technologies/fidelityfx/super-resolution.html}

\end{thebibliography}

\end{document}
